%% main.tex -- "Verified Generators and Relations for Qupit Clifford
%% Operators", five-page version.
%%
%% A compression of ../vqupit (26 pages, OOPSLA/PACMPL format) to five
%% pages of text in the same format, keeping every result, count and
%% design decision of the long version; the rule figure, the normal
%% boxes and the statistics table are in the appendix.  See README.md.
%%
%% Build:  latexmk -pdf main.tex     (or `make'; no Agda run is needed,
%% the code blocks in agda/ are the LaTeX that Agda's backend generated
%% for the long version, copied verbatim).
\documentclass[acmsmall,review,anonymous]{acmart}

%% PACMPL journal metadata (placeholders for the review version).
\acmJournal{PACMPL}
\acmVolume{1}
\acmNumber{OOPSLA}
\acmArticle{1}
\acmMonth{10}
\setcopyright{none}

%% Packages -----------------------------------------------------------
\usepackage{booktabs}
\usepackage{graphicx}
\usepackage{agda-style}      % local: Unicode table for \aid{} in running text
%% Agda's own style file, for the code blocks that `agda --latex' produced
%% from the literate sections of the long version (copied into agda/).
\usepackage{agda}
\renewcommand{\AgdaCodeStyle}{\footnotesize}
\setlength{\mathindent}{1em}     % left indent of code blocks

%% TikZiT circuit figures (generated by the Cir2Tikz tool; the .tikz
%% files are shared with the long version and the QPL 2026 slides) ----
\usepackage{tikz}
\usetikzlibrary{backgrounds,arrows,shapes,shapes.geometric,shapes.misc,cd,arrows.meta}
\pgfdeclarelayer{edgelayer}
\pgfdeclarelayer{nodelayer}
\pgfsetlayers{background,edgelayer,nodelayer,main}
\pgfkeys{/tikz/tikzit fill/.initial=0}
\pgfkeys{/tikz/tikzit draw/.initial=0}
\pgfkeys{/tikz/tikzit shape/.initial=0}
\pgfkeys{/tikz/tikzit category/.initial=0}
\tikzstyle{none}=[inner sep=0mm]
\tikzstyle{tikzfig}=[baseline=-0.25em,scale=0.38]
\newcommand{\tikzfig}[1]{{\tikzstyle{every picture}=[tikzfig]\input{#1.tikz}}}
\definecolor{normGreen}{RGB}{214, 243, 221}
\definecolor{normRed}{RGB}{249, 196, 196}
\definecolor{normYellow}{RGB}{253,246,208}
\definecolor{zxgreen}{RGB}{216,248,216}
\definecolor{zxred}{RGB}{221,165,165}
\definecolor{zxhad}{RGB}{255,255,0}
\definecolor{zxpaulibox}{RGB}{138,159,198}
\input{circuits.tikzstyles}

%% Shrink a circuit to a given width iff it overflows, and vcenter it.
\newlength{\figwidth}
\AtBeginDocument{\setlength{\figwidth}{0.95\textwidth}}
\newsavebox{\figbox}
\newcommand{\fitfig}[2][\figwidth]{%
  \sbox{\figbox}{\tikzfig{figures/#2}}%
  \ifdim\wd\figbox>#1
    \sbox{\figbox}{\resizebox{#1}{!}{\usebox{\figbox}}}%
  \fi
  $\vcenter{\hbox{\usebox{\figbox}}}$}
%% Scale a circuit by a FIXED factor instead: on a grid of relations,
%% every picture shares one scale, so gate boxes stay proportionate.
\newcommand{\ufig}[2][1]{$\vcenter{\hbox{\scalebox{#1}{\tikzfig{figures/#2}}}}$}
%% The same scale applied by TikZ to the coordinates alone, so that the
%% labels of a picture keep their font size (the coset and residual
%% labels b, c, b', c' of the coset-table pictures); the line width is
%% scaled by hand to match the \scalebox'd pictures.
\newcommand{\ufigT}[2][1]{$\vcenter{\hbox{{\tikzstyle{every picture}=[tikzfig,scale=#1,line width=0.4pt*#1]\input{figures/#2.tikz}}}}$}

%% PACMPL reference style (must precede \begin{document})
\citestyle{acmauthoryear}

%% Convenience macros -------------------------------------------------
\newcommand{\pres}[2]{\ensuremath{\langle #1 \mid #2 \rangle}}
\newcommand{\Zp}{\ensuremath{\mathbb{Z}_p}}
\newcommand{\Sp}[1][n]{\ensuremath{\mathrm{Sp}(2#1,\mathbb{Z}_p)}}
\newcommand{\Pauli}[1][n]{\ensuremath{\mathcal{P}_{#1}}}
\newcommand{\Cliff}[1][n]{\ensuremath{\mathcal{C}_{#1}}}
\newcommand{\PCliff}[1][n]{\ensuremath{\overline{\mathcal{C}}_{#1}}}
\newcommand{\PPauli}[1][n]{\ensuremath{\overline{\mathcal{P}}_{#1}}}
\newcommand{\ket}[1]{\ensuremath{|#1\rangle}}
\newcommand{\sem}[1]{\ensuremath{[\![#1]\!]}}
\newcommand{\NF}{\ensuremath{\mathit{NF}}}
\newcommand{\ML}{\ensuremath{\mathit{ML}}}
% Inline Agda identifier (see agda-style.sty for the Unicode table).
\newcommand{\aid}[1]{\mbox{\textsf{\small #1}}}
% Operator / relation variants for math mode, so that $c \aidop{•} c'$
% is set like c • c' rather than c•c'.
\newcommand{\aidop}[1]{\ifmmode\mathbin{\aid{#1}}\else\aid{#1}\fi}
\newcommand{\aidrel}[1]{\ifmmode\mathrel{\aid{#1}}\else\aid{#1}\fi}

\begin{document}

\title{Verified Generators and Relations for Qupit Clifford Operators}

\author{Xiaoning Bian}
\affiliation{%
  \institution{Tsinghua University}
  \city{Beijing}
  \country{China}}
\email{bian@dal.ca}

\author{Yuan Feng}
\affiliation{%
  \institution{Tsinghua University}
  \city{Beijing}
  \country{China}}
\email{yuan\_feng@tsinghua.edu.cn}

\author{Sarah Meng Li}
\affiliation{%
  \institution{University of Amsterdam and QuSoft}
  \city{Amsterdam}
  \country{The Netherlands}}
\email{sarah.li@uva.nl}

\author{John van de Wetering}
\affiliation{%
  \institution{University of Amsterdam and QuSoft}
  \city{Amsterdam}
  \country{The Netherlands}}
\email{john@vdwetering.name}

\begin{abstract}
Bian, Li, Ross, van de Wetering, and Zhao give sixteen rewrite rules
that are sound and complete for multi-qudit Clifford circuits over $H$,
$S$, $CZ$ in every odd prime dimension $p$.  We report a
machine-checked proof of the result in safe, axiom-free Agda, against
an independently constructed semantics.  The proof follows the
paper's for the symplectic factor, where the paper does its real work,
and diverges once, at the step that promotes symplectic completeness
to projective completeness ``up to Paulis'': there we factor the projective
Clifford group as $\overline{\mathcal P}_n \rtimes
\mathrm{Sp}(2n,\mathbb{Z}_p)$, present the factors separately, and glue
the presentations with a generic, verified semidirect-product theorem.
The symplectic factor carries the mathematics: the paper's normal form,
refined into a doubly inductive datatype of \emph{normal boxes} and
\emph{rows}, is a tower of verified Reidemeister--Schreier coset
tables, proved unique by induction on the width, and the paper's 42 box
relations are derived from eighteen natural ones.  Verified
isomorphisms of presentations reach the paper's own rules; a central
extension, verified against a cocycle, and a factor $\langle -1\rangle$
restore the scalars, reaching Figure~1 as printed, over $-\omega$.  Everything is parametric in $p$ and typechecks with no postulates.

\end{abstract}

%% CCS concepts (required by acmart) ----------------------------------
\begin{CCSXML}
<ccs2012>
<concept>
<concept_id>10003752.10003790.10011740</concept_id>
<concept_desc>Theory of computation~Type theory</concept_desc>
<concept_significance>500</concept_significance>
</concept>
<concept>
<concept_id>10003752.10003790.10002990</concept_id>
<concept_desc>Theory of computation~Logic and verification</concept_desc>
<concept_significance>500</concept_significance>
</concept>
<concept>
<concept_id>10003752.10003790.10003800</concept_id>
<concept_desc>Theory of computation~Equational logic and rewriting</concept_desc>
<concept_significance>500</concept_significance>
</concept>
<concept>
<concept_id>10003752.10010124.10010131.10010137</concept_id>
<concept_desc>Theory of computation~Quantum computation theory</concept_desc>
<concept_significance>500</concept_significance>
</concept>
</ccs2012>
\end{CCSXML}
\ccsdesc[500]{Theory of computation~Type theory}
\ccsdesc[500]{Theory of computation~Logic and verification}
\ccsdesc[500]{Theory of computation~Equational logic and rewriting}
\ccsdesc[500]{Theory of computation~Quantum computation theory}

\keywords{Clifford group, qudits, generators and relations, completeness,
normal forms, symplectic group, semidirect product, Agda, formal
verification}

\maketitle

\section{Introduction}\label{sec:intro}

A finite list of circuit equations is \emph{sound} when each holds in
the intended semantics and \emph{complete} when every semantic equality
between circuits is derivable from it.  Completeness is proved by a
\emph{normal form}, and the case analysis of every generator against
every normal form is enormous:
Selinger's $n$-qubit Clifford presentation~\cite{selinger2015clifford}
runs a substantial rewrite system through it, its qutrit
successor~\cite{li2025qutrit} needs over fifty pages, and its
generalisation to all odd primes by Bian, Li, Ross, van de Wetering, and
Zhao~\cite{bian2026qudit} --- sixteen rules for multi-qudit Clifford
circuits in every odd prime dimension $p$, henceforth ``the qupit
paper'' --- rests on roughly a hundred pages of
derivations~\cite{bian2026supplement}.

We report a machine-checked proof of that theorem in safe, axiom-free
Agda~\cite{norell2007thesis,agdastdlib}, against a semantics constructed
independently of the syntax, on the library of~\cite{bian2026framework,qupit-code} (introduced in Appendix~\ref{sec:permutations}).  The
proof follows the paper's for the symplectic factor, where the paper
does its real work, and diverges once, at the step that promotes
symplectic to projective completeness \emph{up to Paulis}
(\S\ref{sec:background}); there we instead \emph{factor the group,
present the factors, compose the presentations}: $\PCliff = \PPauli
\rtimes \Sp$ and $\Cliff$ is a central extension of it by the scalars;
the symplectic group gets a coset-table normal form, the Pauli group a
compositional presentation, the scalars a cyclic one, generic verified
theorems for semidirect products and central extensions glue them,
plumbing reduces the result to the paper's $H, S, CZ$ and rules, and a
factor $\langle -1\rangle$ reaches its Figure~1 over $-\omega$.

\section{Qupit Clifford Circuits and the Main Theorems}\label{sec:background}

Fix an odd prime $p$ and $\omega = e^{2\pi i/p}$.  On the basis
$\ket{0},\dots,\ket{p-1}$ indexed by $\Zp$, $X\ket{j} = \ket{j+1}$ and
$Z\ket{j} = \omega^{j}\ket{j}$ generate, with the phases $\omega^t$, the
Pauli group $\Pauli$, and $H\colon \ket{j} \mapsto
\tfrac{\delta_p}{\sqrt p}\sum_k \omega^{jk}\ket{k}$, $S\colon \ket{j}
\mapsto \omega^{j(j-1)/2}\ket{j}$, $CZ\colon \ket{j}\ket{l} \mapsto
\omega^{jl}\ket{j}\ket{l}$ generate, with the phases $\omega^t$, the
Clifford group $\Cliff$ ($\delta_p = i^{(p-1)/2}$ makes the
determinant of $H$ equal to $1$; the determinants of $S$ and $CZ$ are
$1$ too)~\cite{bian2026qudit}.  Modulo scalars, $\PCliff =
\Cliff/\langle\omega\rangle$ and $\PPauli \cong \Zp^{2n}$, $\Cliff/\Pauli \cong \PCliff/\PPauli \cong \Sp$, and $\PCliff \cong
\PPauli \rtimes \Sp$.  These ``glue facts'' are the two not formalised:
our semantic groups are \emph{defined} as the right-hand sides.
In~\cite{bian2026qudit} the global phases also include $-1$, so that
the swap and the multiplier gates belong to $\Cliff$; we work mainly
in the projective quotient, where $-1$ makes no difference except that
the relations come out a little simpler.

Using the framework of Appendix~\ref{sec:permutations}, we encode
circuits, circuit relations, and reasoning with the relations.  We
draw circuits in matrix-multiplication order, the reverse of the usual
circuit order: $w \bullet v$ applies $v$ first but draws it on the
right.  For example, the swap:

\noindent\begin{minipage}[c]{0.56\textwidth}
%% Generated by make agda from lagda/ex-circuit.lagda.tex (agda --latex --only-scope-checking);
%% edit that file, not this one.  The hidden block renders nothing.
\begin{code}[hide]%
\>[0]\AgdaKeyword{module}\AgdaSpace{}%
\AgdaModule{vqupit-short.lagda.ex-circuit}\AgdaSpace{}%
\AgdaKeyword{where}\<%
\\
%
\\[\AgdaEmptyExtraSkip]%
\>[0]\AgdaKeyword{open}\AgdaSpace{}%
\AgdaKeyword{import}\AgdaSpace{}%
\AgdaModule{vqupit-short.lagda.Prelude}\<%
\\
%
\\[\AgdaEmptyExtraSkip]%
\>[0]\AgdaKeyword{private}\AgdaSpace{}%
\AgdaKeyword{variable}\<%
\\
\>[0][@{}l@{\AgdaIndent{0}}]%
\>[2]\AgdaGeneralizable{n}\AgdaSpace{}%
\AgdaSymbol{:}\AgdaSpace{}%
\AgdaDatatype{ℕ}\<%
\end{code}
\begin{code}%
\>[0]\AgdaFunction{Ex}\AgdaSpace{}%
\AgdaSymbol{:}\AgdaSpace{}%
\AgdaDatatype{Word}\AgdaSpace{}%
\AgdaSymbol{(}\AgdaPostulate{Gen}\AgdaSpace{}%
\AgdaSymbol{(}\AgdaInductiveConstructor{₂₊}\AgdaSpace{}%
\AgdaGeneralizable{n}\AgdaSymbol{))}\<%
\\
\>[0]\AgdaFunction{Ex}\AgdaSpace{}%
\AgdaSymbol{=}\AgdaSpace{}%
\AgdaPostulate{CZ}\AgdaSpace{}%
\AgdaOperator{\AgdaInductiveConstructor{•}}\AgdaSpace{}%
\AgdaPostulate{H}\AgdaSpace{}%
\AgdaOperator{\AgdaFunction{↓}}\AgdaSpace{}%
\AgdaOperator{\AgdaInductiveConstructor{•}}\AgdaSpace{}%
\AgdaPostulate{H}\AgdaSpace{}%
\AgdaOperator{\AgdaPostulate{↑}}\AgdaSpace{}%
\AgdaOperator{\AgdaInductiveConstructor{•}}\AgdaSpace{}%
\AgdaPostulate{CZ}\AgdaSpace{}%
\AgdaOperator{\AgdaInductiveConstructor{•}}\AgdaSpace{}%
\AgdaPostulate{H}\AgdaSpace{}%
\AgdaOperator{\AgdaFunction{↓}}\AgdaSpace{}%
\AgdaOperator{\AgdaInductiveConstructor{•}}\AgdaSpace{}%
\AgdaPostulate{H}\AgdaSpace{}%
\AgdaOperator{\AgdaPostulate{↑}}\AgdaSpace{}%
\AgdaOperator{\AgdaInductiveConstructor{•}}\AgdaSpace{}%
\AgdaPostulate{CZ}\AgdaSpace{}%
\AgdaOperator{\AgdaInductiveConstructor{•}}\AgdaSpace{}%
\AgdaPostulate{H}\AgdaSpace{}%
\AgdaOperator{\AgdaFunction{↓}}\AgdaSpace{}%
\AgdaOperator{\AgdaInductiveConstructor{•}}\AgdaSpace{}%
\AgdaPostulate{H}\AgdaSpace{}%
\AgdaOperator{\AgdaPostulate{↑}}\<%
\end{code}

\end{minipage}\hfill
\begin{minipage}[c]{0.42\textwidth}\centering
\ufig[0.48]{def-Ex-eq}
\end{minipage}

\noindent A rule set is an inductive family indexed by the width, one
constructor per rule, each stated once at the smallest width at which
it fits; for example, two of the fifteen constructors of our primary
rule set \aid{{\_}SRel,{\_}==={\_}}, the order of the swap and the
rule moving $H$ across it:

\noindent\begin{minipage}[c]{0.56\textwidth}
%% Generated by make agda from lagda/ex-relation.lagda.tex (agda --latex --only-scope-checking);
%% edit that file, not this one.  The hidden block renders nothing.
\begin{code}[hide]%
\>[0]\AgdaKeyword{module}\AgdaSpace{}%
\AgdaModule{vqupit-short.lagda.ex-relation}\AgdaSpace{}%
\AgdaKeyword{where}\<%
\\
%
\\[\AgdaEmptyExtraSkip]%
\>[0]\AgdaKeyword{open}\AgdaSpace{}%
\AgdaKeyword{import}\AgdaSpace{}%
\AgdaModule{vqupit-short.lagda.Prelude}\<%
\\
\>[0]\AgdaKeyword{open}\AgdaSpace{}%
\AgdaKeyword{import}\AgdaSpace{}%
\AgdaModule{vqupit-short.lagda.ex-circuit}\<%
\\
%
\\[\AgdaEmptyExtraSkip]%
\>[0]\AgdaKeyword{private}\AgdaSpace{}%
\AgdaKeyword{variable}\<%
\\
\>[0][@{}l@{\AgdaIndent{0}}]%
\>[2]\AgdaGeneralizable{n}\AgdaSpace{}%
\AgdaSymbol{:}\AgdaSpace{}%
\AgdaDatatype{ℕ}\<%
\\
%
\\[\AgdaEmptyExtraSkip]%
\>[0]\AgdaKeyword{infix}\AgdaSpace{}%
\AgdaNumber{4}\AgdaSpace{}%
\AgdaOperator{\AgdaDatatype{\AgdaUnderscore{}CRel,\AgdaUnderscore{}===\AgdaUnderscore{}}}\<%
\end{code}
\begin{code}%
\>[0]\AgdaKeyword{data}\AgdaSpace{}%
\AgdaOperator{\AgdaDatatype{\AgdaUnderscore{}CRel,\AgdaUnderscore{}===\AgdaUnderscore{}}}\AgdaSpace{}%
\AgdaSymbol{:}\AgdaSpace{}%
\AgdaSymbol{(}\AgdaBound{n}\AgdaSpace{}%
\AgdaSymbol{:}\AgdaSpace{}%
\AgdaDatatype{ℕ}\AgdaSymbol{)}\AgdaSpace{}%
\AgdaSymbol{→}\AgdaSpace{}%
\AgdaFunction{WRel}\AgdaSpace{}%
\AgdaSymbol{(}\AgdaPostulate{Gen}\AgdaSpace{}%
\AgdaBound{n}\AgdaSymbol{)}\AgdaSpace{}%
\AgdaKeyword{where}\<%
\\
\>[0][@{}l@{\AgdaIndent{0}}]%
\>[2]\AgdaInductiveConstructor{order-Ex}%
\>[14]\AgdaSymbol{:}\AgdaSpace{}%
\AgdaSymbol{(}\AgdaInductiveConstructor{₂₊}\AgdaSpace{}%
\AgdaGeneralizable{n}\AgdaSymbol{)}\AgdaSpace{}%
\AgdaOperator{\AgdaDatatype{CRel,}}%
\>[30]\AgdaFunction{Ex}\AgdaSpace{}%
\AgdaOperator{\AgdaFunction{\textasciicircum{}}}\AgdaSpace{}%
\AgdaNumber{2}\AgdaSpace{}%
\AgdaOperator{\AgdaDatatype{===}}\AgdaSpace{}%
\AgdaInductiveConstructor{ε}\<%
\\
%
\>[2]\AgdaInductiveConstructor{semi-Ex-H↑}%
\>[14]\AgdaSymbol{:}\AgdaSpace{}%
\AgdaSymbol{(}\AgdaInductiveConstructor{₂₊}\AgdaSpace{}%
\AgdaGeneralizable{n}\AgdaSymbol{)}\AgdaSpace{}%
\AgdaOperator{\AgdaDatatype{CRel,}}%
\>[30]\AgdaFunction{Ex}\AgdaSpace{}%
\AgdaOperator{\AgdaInductiveConstructor{•}}\AgdaSpace{}%
\AgdaPostulate{H}\AgdaSpace{}%
\AgdaOperator{\AgdaPostulate{↑}}\AgdaSpace{}%
\AgdaOperator{\AgdaDatatype{===}}\AgdaSpace{}%
\AgdaPostulate{H}\AgdaSpace{}%
\AgdaOperator{\AgdaFunction{↓}}\AgdaSpace{}%
\AgdaOperator{\AgdaInductiveConstructor{•}}\AgdaSpace{}%
\AgdaFunction{Ex}\<%
\end{code}

\end{minipage}\hfill
\begin{minipage}[c]{0.42\textwidth}\centering
\ufig[0.48]{pv0-order-Ex}\\[3pt]
\ufig[0.48]{pv0-semi-Ex-Hup}
\end{minipage}

\noindent And reasoning with the rules is equational reasoning in the
congruence $\approx$ they generate, with combinators for symmetry,
associativity, and rewriting the left or the right factor of a
product; for example, conjugating by the swap moves $H$ down a wire:

\noindent\begin{minipage}[c]{0.56\textwidth}
%% Generated by make agda from lagda/ex-reasoning.lagda.tex (agda --latex --only-scope-checking);
%% edit that file, not this one.  The hidden block renders nothing.
\begin{code}[hide]%
\>[0]\AgdaKeyword{module}\AgdaSpace{}%
\AgdaModule{vqupit-short.lagda.ex-reasoning}\AgdaSpace{}%
\AgdaKeyword{where}\<%
\\
%
\\[\AgdaEmptyExtraSkip]%
\>[0]\AgdaKeyword{open}\AgdaSpace{}%
\AgdaKeyword{import}\AgdaSpace{}%
\AgdaModule{vqupit-short.lagda.Prelude}\<%
\\
\>[0]\AgdaKeyword{open}\AgdaSpace{}%
\AgdaKeyword{import}\AgdaSpace{}%
\AgdaModule{vqupit-short.lagda.ex-circuit}\<%
\\
\>[0]\AgdaKeyword{open}\AgdaSpace{}%
\AgdaKeyword{import}\AgdaSpace{}%
\AgdaModule{vqupit-short.lagda.ex-relation}\<%
\\
%
\\[\AgdaEmptyExtraSkip]%
\>[0]\AgdaKeyword{private}\AgdaSpace{}%
\AgdaKeyword{variable}\<%
\\
\>[0][@{}l@{\AgdaIndent{0}}]%
\>[2]\AgdaGeneralizable{n}\AgdaSpace{}%
\AgdaSymbol{:}\AgdaSpace{}%
\AgdaDatatype{ℕ}\<%
\\
%
\\[\AgdaEmptyExtraSkip]%
\>[0]\AgdaKeyword{postulate}\<%
\\
\>[0][@{}l@{\AgdaIndent{0}}]%
\>[2]\AgdaPostulate{axiom}\AgdaSpace{}%
\AgdaSymbol{:}\AgdaSpace{}%
\AgdaSymbol{∀}\AgdaSpace{}%
\AgdaSymbol{\{}\AgdaBound{w}\AgdaSpace{}%
\AgdaBound{v}\AgdaSpace{}%
\AgdaSymbol{:}\AgdaSpace{}%
\AgdaDatatype{Word}\AgdaSpace{}%
\AgdaSymbol{(}\AgdaPostulate{Gen}\AgdaSpace{}%
\AgdaGeneralizable{n}\AgdaSymbol{)\}}\AgdaSpace{}%
\AgdaSymbol{→}\AgdaSpace{}%
\AgdaGeneralizable{n}\AgdaSpace{}%
\AgdaOperator{\AgdaDatatype{SRel,}}\AgdaSpace{}%
\AgdaBound{w}\AgdaSpace{}%
\AgdaOperator{\AgdaDatatype{===}}\AgdaSpace{}%
\AgdaBound{v}\AgdaSpace{}%
\AgdaSymbol{→}\AgdaSpace{}%
\AgdaBound{w}\AgdaSpace{}%
\AgdaOperator{\AgdaPostulate{≈}}\AgdaSpace{}%
\AgdaBound{v}\<%
\end{code}
\begin{code}%
\>[0]\AgdaFunction{conj-Ex-H↑}\AgdaSpace{}%
\AgdaSymbol{:}\AgdaSpace{}%
\AgdaFunction{Ex}\AgdaSpace{}%
\AgdaOperator{\AgdaInductiveConstructor{•}}\AgdaSpace{}%
\AgdaPostulate{H}\AgdaSpace{}%
\AgdaOperator{\AgdaPostulate{↑}}\AgdaSpace{}%
\AgdaOperator{\AgdaInductiveConstructor{•}}\AgdaSpace{}%
\AgdaFunction{Ex}\AgdaSpace{}%
\AgdaOperator{\AgdaPostulate{≈}}\AgdaSpace{}%
\AgdaPostulate{H}\AgdaSpace{}%
\AgdaOperator{\AgdaFunction{↓}}\<%
\\
\>[0]\AgdaFunction{conj-Ex-H↑}\AgdaSpace{}%
\AgdaSymbol{=}\AgdaSpace{}%
\AgdaOperator{\AgdaPostulate{begin}}\<%
\\
\>[0][@{}l@{\AgdaIndent{0}}]%
\>[2]\AgdaFunction{Ex}\AgdaSpace{}%
\AgdaOperator{\AgdaInductiveConstructor{•}}\AgdaSpace{}%
\AgdaPostulate{H}\AgdaSpace{}%
\AgdaOperator{\AgdaPostulate{↑}}\AgdaSpace{}%
\AgdaOperator{\AgdaInductiveConstructor{•}}\AgdaSpace{}%
\AgdaFunction{Ex}%
\>[19]\AgdaOperator{\AgdaPostulate{≈⟨}}\AgdaSpace{}%
\AgdaPostulate{sym}\AgdaSpace{}%
\AgdaPostulate{assoc}\AgdaSpace{}%
\AgdaOperator{\AgdaPostulate{⟩}}\<%
\\
%
\>[2]\AgdaSymbol{(}\AgdaFunction{Ex}\AgdaSpace{}%
\AgdaOperator{\AgdaInductiveConstructor{•}}\AgdaSpace{}%
\AgdaPostulate{H}\AgdaSpace{}%
\AgdaOperator{\AgdaPostulate{↑}}\AgdaSymbol{)}\AgdaSpace{}%
\AgdaOperator{\AgdaInductiveConstructor{•}}\AgdaSpace{}%
\AgdaFunction{Ex}%
\>[19]\AgdaOperator{\AgdaPostulate{≈⟨}}\AgdaSpace{}%
\AgdaOperator{\AgdaPostulate{cleft}}\AgdaSpace{}%
\AgdaPostulate{axiom}\AgdaSpace{}%
\AgdaInductiveConstructor{semi-Ex-H↑}\AgdaSpace{}%
\AgdaOperator{\AgdaPostulate{⟩}}\<%
\\
%
\>[2]\AgdaSymbol{(}\AgdaPostulate{H}\AgdaSpace{}%
\AgdaOperator{\AgdaFunction{↓}}\AgdaSpace{}%
\AgdaOperator{\AgdaInductiveConstructor{•}}\AgdaSpace{}%
\AgdaFunction{Ex}\AgdaSymbol{)}\AgdaSpace{}%
\AgdaOperator{\AgdaInductiveConstructor{•}}\AgdaSpace{}%
\AgdaFunction{Ex}%
\>[19]\AgdaOperator{\AgdaPostulate{≈⟨}}\AgdaSpace{}%
\AgdaPostulate{assoc}\AgdaSpace{}%
\AgdaOperator{\AgdaPostulate{⟩}}\<%
\\
%
\>[2]\AgdaPostulate{H}\AgdaSpace{}%
\AgdaOperator{\AgdaFunction{↓}}\AgdaSpace{}%
\AgdaOperator{\AgdaInductiveConstructor{•}}\AgdaSpace{}%
\AgdaFunction{Ex}\AgdaSpace{}%
\AgdaOperator{\AgdaInductiveConstructor{•}}\AgdaSpace{}%
\AgdaFunction{Ex}%
\>[19]\AgdaOperator{\AgdaPostulate{≈⟨}}\AgdaSpace{}%
\AgdaOperator{\AgdaPostulate{cright}}\AgdaSpace{}%
\AgdaPostulate{axiom}\AgdaSpace{}%
\AgdaInductiveConstructor{order-Ex}\AgdaSpace{}%
\AgdaOperator{\AgdaPostulate{⟩}}\<%
\\
%
\>[2]\AgdaPostulate{H}\AgdaSpace{}%
\AgdaOperator{\AgdaFunction{↓}}\AgdaSpace{}%
\AgdaOperator{\AgdaInductiveConstructor{•}}\AgdaSpace{}%
\AgdaInductiveConstructor{ε}%
\>[19]\AgdaOperator{\AgdaPostulate{≈⟨}}\AgdaSpace{}%
\AgdaPostulate{right-unit}\AgdaSpace{}%
\AgdaOperator{\AgdaPostulate{⟩}}\<%
\\
%
\>[2]\AgdaPostulate{H}\AgdaSpace{}%
\AgdaOperator{\AgdaFunction{↓}}\AgdaSpace{}%
\AgdaOperator{\AgdaPostulate{∎}}\<%
\end{code}

\end{minipage}\hfill
\begin{minipage}[c]{0.42\textwidth}\centering
\ufig[0.48]{pv0-conj-Ex-Hup-steps}
\end{minipage}

\noindent Further, we can express and prove our main theorems:

\begin{theorem}[\aid{clifford-presentation}]\label{thm:main}
For every width $n$, the fifteen rules---the relations
of~\cite{bian2026qudit}, read modulo scalars---present the projective
Clifford group:
%% Copied from ../vqupit/latex/vqupit/sections/composing.tex (composing, the clifford-presentation block), the LaTeX
%% that `agda --latex --only-scope-checking' generated for the long version.
\begin{code}%
\>[0]\AgdaFunction{clifford-presentation}\AgdaSpace{}%
\AgdaSymbol{:}\AgdaSpace{}%
\AgdaSymbol{∀}\AgdaSpace{}%
\AgdaBound{n}\AgdaSpace{}%
\AgdaSymbol{→}\AgdaSpace{}%
\AgdaSymbol{(}\AgdaOperator{\AgdaPostulate{PapV1.Clifford-Relations.\AgdaUnderscore{}QRel,\AgdaUnderscore{}===\AgdaUnderscore{}}}\AgdaSpace{}%
\AgdaBound{n}\AgdaSymbol{)}\AgdaSpace{}%
\AgdaOperator{\AgdaRecord{IsPresentationOf}}\AgdaSpace{}%
\AgdaSymbol{(}\AgdaPostulate{Pauli⋊Sp}\AgdaSpace{}%
\AgdaBound{n}\AgdaSymbol{)}\<%
\end{code}

\end{theorem}

\noindent \aid{{\_}QRel,{\_}==={\_}} closes \aid{{\_}SRel,{\_}==={\_}}
under the three structural rules (Appendix~\ref{sec:permutations});
$\aid{Pauli⋊Sp}\;n$ is the semidirect product of $\Zp^{2n}$
by $\Sp$ (\S\ref{sec:composing}); \aid{{\_}IsPresentationOf{\_}}
(Appendix~\ref{sec:permutations}) is a group isomorphism $\mathsf{Cir}\;n/{\approx}
\to \PPauli \rtimes \Sp$:
\emph{soundness} (well-definedness of the map), each rule holds;
\emph{completeness}, $\sem{\_}$ is injective; and \emph{surjectivity},
each projective Clifford operator is represented by a circuit.

We refine the qupit paper's normal form for symplectic operators
(Clifford operators modulo Paulis) into a doubly inductive definition
(\S\ref{sec:symplectic}) and prove that it
enjoys the unique-normal-form property:

\begin{theorem}[\aid{unique-nf}]\label{thm:unique}
Normal forms with equal symplectic denotations are equal:
%% Copied from ../vqupit/latex/vqupit/sections/background.tex (background, block 4), the LaTeX
%% that `agda --latex --only-scope-checking' generated for the long version.
\begin{code}%
\>[0]\AgdaFunction{Theorem-2}\AgdaSpace{}%
\AgdaSymbol{:}\AgdaSpace{}%
\AgdaSymbol{∀}\AgdaSpace{}%
\AgdaBound{n}\AgdaSpace{}%
\AgdaSymbol{\{}\AgdaBound{u}\AgdaSpace{}%
\AgdaBound{v}\AgdaSpace{}%
\AgdaSymbol{:}\AgdaSpace{}%
\AgdaPostulate{NF}\AgdaSpace{}%
\AgdaBound{n}\AgdaSymbol{\}}\AgdaSpace{}%
\AgdaSymbol{→}\AgdaSpace{}%
\AgdaOperator{\AgdaPostulate{⟦}}\AgdaSpace{}%
\AgdaOperator{\AgdaPostulate{[}}\AgdaSpace{}%
\AgdaBound{u}\AgdaSpace{}%
\AgdaOperator{\AgdaPostulate{]}}\AgdaSpace{}%
\AgdaOperator{\AgdaPostulate{⟧}}\AgdaSpace{}%
\AgdaOperator{\AgdaPostulate{≈ˢ}}\AgdaSpace{}%
\AgdaOperator{\AgdaPostulate{⟦}}\AgdaSpace{}%
\AgdaOperator{\AgdaPostulate{[}}\AgdaSpace{}%
\AgdaBound{v}\AgdaSpace{}%
\AgdaOperator{\AgdaPostulate{]}}\AgdaSpace{}%
\AgdaOperator{\AgdaPostulate{⟧}}\AgdaSpace{}%
\AgdaSymbol{→}\AgdaSpace{}%
\AgdaBound{u}\AgdaSpace{}%
\AgdaOperator{\AgdaDatatype{≡}}\AgdaSpace{}%
\AgdaBound{v}\<%
\\
\>[0]\AgdaFunction{Theorem-2}\AgdaSpace{}%
\AgdaSymbol{=}\AgdaSpace{}%
\AgdaPostulate{U.⟦[]⟧-injective}\<%
\end{code}

\end{theorem}

\noindent\emph{The route.}  We mostly follow the paper: its normal
boxes, its symplectic normal form and its
parametrised \emph{box relations} --- pushing a gate through a normal
box yields an updated box and residual \emph{dirty} gates --- reappear
below, with Theorem~4.4's reduction to eighteen relations, a step the
paper already delegated to Agda.  The one major divergence is the step
that promotes symplectic to projective completeness (its
Proposition~4.9): the paper corrects each rule $C_l = C_r$ to $C_l =
C_r\,;P$ for the unique Pauli $P$ making it projectively sound and
replays the symplectic rewrites, pushing each correction to the end,
where it becomes ``a potentially different Pauli''.  We did not see
how to verify it at reasonable cost.  We instead work in a syntax where Paulis are
genuinely trivial --- symplectic circuits have only $H$, $S$, $CZ$ and
denote $\Sp$ itself --- and reintroduce Paulis and scalars by
construction (Figure~\ref{fig:route} in the appendix): the Pauli
bookkeeping of Proposition~4.9 is done once, by a semidirect-product
theorem proved for arbitrary groups before this project began.

\section{The Symplectic Factor}\label{sec:symplectic}

$\Sp$ is modelled as the group of $\Zp$-linear, form-preserving
bijections of the phase space $\aid{Pauli}\;n = \aid{Vec}\;(\Zp \times
\Zp)\;n$ (the exponents of $Z^{a}X^{b}$ up to phase, per wire).
Against this semantics, completeness rests on the normal form, which
we refine into a doubly inductive datatype:

\noindent\begin{minipage}[c]{0.4\textwidth}
\begin{code}%
\>[0]\AgdaFunction{M}\AgdaSpace{}%
\AgdaFunction{L}\AgdaSpace{}%
\AgdaFunction{ML}\AgdaSpace{}%
\AgdaSymbol{:}\AgdaSpace{}%
\AgdaDatatype{ℕ}\AgdaSpace{}%
\AgdaSymbol{→}\AgdaSpace{}%
\AgdaPrimitive{Set}\<%
\\
\>[0]\AgdaFunction{M}\AgdaSpace{}%
\AgdaNumber{0}\AgdaSpace{}%
\AgdaSymbol{=}\AgdaSpace{}%
\AgdaRecord{⊤}\<%
\\
\>[0]\AgdaFunction{M}\AgdaSpace{}%
\AgdaSymbol{(}\AgdaInductiveConstructor{₁₊}\AgdaSpace{}%
\AgdaBound{n}\AgdaSymbol{)}\AgdaSpace{}%
\AgdaSymbol{=}\AgdaSpace{}%
\AgdaDatatype{Vec}\AgdaSpace{}%
\AgdaPostulate{D}\AgdaSpace{}%
\AgdaBound{n}\AgdaSpace{}%
\AgdaOperator{\AgdaFunction{×}}\AgdaSpace{}%
\AgdaPostulate{E}\<%
\\
\>[0]\AgdaFunction{L}\AgdaSpace{}%
\AgdaNumber{0}\AgdaSpace{}%
\AgdaSymbol{=}\AgdaSpace{}%
\AgdaRecord{⊤}\<%
\\
\>[0]\AgdaFunction{L}\AgdaSpace{}%
\AgdaSymbol{(}\AgdaInductiveConstructor{₁₊}\AgdaSpace{}%
\AgdaBound{n}\AgdaSymbol{)}\AgdaSpace{}%
\AgdaSymbol{=}\AgdaSpace{}%
\AgdaDatatype{Vec}\AgdaSpace{}%
\AgdaPostulate{B}\AgdaSpace{}%
\AgdaBound{n}\AgdaSpace{}%
\AgdaOperator{\AgdaFunction{×}}\AgdaSpace{}%
\AgdaPostulate{A}\<%
\\
\>[0]\AgdaFunction{ML}\AgdaSpace{}%
\AgdaNumber{0}\AgdaSpace{}%
\AgdaSymbol{=}\AgdaSpace{}%
\AgdaRecord{⊤}\<%
\\
\>[0]\AgdaFunction{ML}\AgdaSpace{}%
\AgdaNumber{1}\AgdaSpace{}%
\AgdaSymbol{=}\AgdaSpace{}%
\AgdaFunction{M}\AgdaSpace{}%
\AgdaNumber{1}\AgdaSpace{}%
\AgdaOperator{\AgdaFunction{×}}\AgdaSpace{}%
\AgdaFunction{L}\AgdaSpace{}%
\AgdaNumber{1}\<%
\\
\>[0]\AgdaFunction{ML}\AgdaSpace{}%
\AgdaBound{m}\AgdaSymbol{@(}\AgdaInductiveConstructor{₂₊}\AgdaSpace{}%
\AgdaBound{n}\AgdaSymbol{)}%
\>[350I]\AgdaSymbol{=}\AgdaSpace{}%
\AgdaFunction{M}\AgdaSpace{}%
\AgdaBound{m}\AgdaSpace{}%
\AgdaOperator{\AgdaFunction{×}}\AgdaSpace{}%
\AgdaFunction{L}\AgdaSpace{}%
\AgdaBound{m}\<%
\\
\>[.][@{}l@{}]\<[350I]%
\>[12]\AgdaOperator{\AgdaDatatype{⊎}}\AgdaSpace{}%
\AgdaPostulate{D}\AgdaSpace{}%
\AgdaOperator{\AgdaFunction{×}}\AgdaSpace{}%
\AgdaFunction{ML}\AgdaSpace{}%
\AgdaSymbol{(}\AgdaInductiveConstructor{₁₊}\AgdaSpace{}%
\AgdaBound{n}\AgdaSymbol{)}\<%
\end{code}

\end{minipage}\hfill
\begin{minipage}[c]{0.56\textwidth}\centering
  \begin{tabular}{@{}r@{\,}c@{\,}l@{\quad}r@{\,}c@{\,}l@{}}
    \ufig[0.5]{m4-box} & \scalebox{0.7}{$=$} & \ufig[0.5]{m4-atoms} &
    \ufig[0.5]{ml4-box} & \scalebox{0.7}{$=$} & \ufig[0.5]{ml-tower-a-r}\\[1ex]
    \ufig[0.5]{l4-box} & \scalebox{0.7}{$=$} & \ufig[0.5]{l4-atoms} &
    \ufig[0.5]{ml4-box} & \scalebox{0.7}{$=$} & \ufig[0.5]{ml-tower-b-r}\\
  \end{tabular}
\end{minipage}

\medskip\noindent
An $M$-row is a column of $D$ boxes ending in an $E$ box (the paper's
$X$-normal layer); an $L$-row is a column of $B$ boxes ending in an $A$
box (its $Z$-normal layer, with the $A$ box at the bottom).  An
$\ML$ box at width $m \ge 2$ is \emph{either} a full pair of rows
\emph{or} a bottom $D$ box followed by an $\ML$ box one wire up: a
spine of $D$ boxes capped by a pair of rows.  This is the first
induction, on the width of a single coset representative.  The second
induction is the staircase:

\noindent\begin{minipage}[c]{0.54\textwidth}
\begin{code}%
\>[0]\AgdaFunction{NF}\AgdaSpace{}%
\AgdaSymbol{:}\AgdaSpace{}%
\AgdaDatatype{ℕ}\AgdaSpace{}%
\AgdaSymbol{→}\AgdaSpace{}%
\AgdaPrimitive{Set}\<%
\\
\>[0]\AgdaFunction{NF}\AgdaSpace{}%
\AgdaNumber{0}\AgdaSpace{}%
\AgdaSymbol{=}\AgdaSpace{}%
\AgdaRecord{⊤}\<%
\\
\>[0]\AgdaFunction{NF}\AgdaSpace{}%
\AgdaSymbol{(}\AgdaInductiveConstructor{₁₊}\AgdaSpace{}%
\AgdaBound{n}\AgdaSymbol{)}\AgdaSpace{}%
\AgdaSymbol{=}\AgdaSpace{}%
\AgdaFunction{NF}\AgdaSpace{}%
\AgdaBound{n}\AgdaSpace{}%
\AgdaOperator{\AgdaFunction{×}}\AgdaSpace{}%
\AgdaFunction{ML}\AgdaSpace{}%
\AgdaSymbol{(}\AgdaInductiveConstructor{₁₊}\AgdaSpace{}%
\AgdaBound{n}\AgdaSymbol{)}\<%
\\
%
\\[\AgdaEmptyExtraSkip]%
\>[0]\AgdaOperator{\AgdaFunction{[\AgdaUnderscore{}]}}\AgdaSpace{}%
\AgdaSymbol{:}\AgdaSpace{}%
\AgdaSymbol{∀}\AgdaSpace{}%
\AgdaSymbol{\{}\AgdaBound{n}\AgdaSymbol{\}}\AgdaSpace{}%
\AgdaSymbol{→}\AgdaSpace{}%
\AgdaFunction{NF}\AgdaSpace{}%
\AgdaBound{n}\AgdaSpace{}%
\AgdaSymbol{→}\AgdaSpace{}%
\AgdaDatatype{Word}\AgdaSpace{}%
\AgdaSymbol{(}\AgdaDatatype{Gen}\AgdaSpace{}%
\AgdaBound{n}\AgdaSymbol{)}\<%
\\
\>[0]\AgdaOperator{\AgdaFunction{[\AgdaUnderscore{}]}}\AgdaSpace{}%
\AgdaSymbol{\{}\AgdaNumber{0}\AgdaSymbol{\}}%
\>[12]\AgdaInductiveConstructor{tt}%
\>[23]\AgdaSymbol{=}\AgdaSpace{}%
\AgdaInductiveConstructor{ε}\<%
\\
\>[0]\AgdaOperator{\AgdaFunction{[\AgdaUnderscore{}]}}\AgdaSpace{}%
\AgdaSymbol{\{}\AgdaInductiveConstructor{₁₊}\AgdaSpace{}%
\AgdaBound{n}\AgdaSymbol{\}}%
\>[12]\AgdaSymbol{(}\AgdaBound{nf}\AgdaSpace{}%
\AgdaOperator{\AgdaInductiveConstructor{,}}\AgdaSpace{}%
\AgdaBound{ml}\AgdaSymbol{)}%
\>[23]\AgdaSymbol{=}\AgdaSpace{}%
\AgdaOperator{\AgdaFunction{[}}\AgdaSpace{}%
\AgdaBound{nf}\AgdaSpace{}%
\AgdaOperator{\AgdaFunction{]}}\AgdaSpace{}%
\AgdaOperator{\AgdaPostulate{↑}}\AgdaSpace{}%
\AgdaOperator{\AgdaInductiveConstructor{•}}\AgdaSpace{}%
\AgdaOperator{\AgdaPostulate{[}}\AgdaSpace{}%
\AgdaBound{ml}\AgdaSpace{}%
\AgdaOperator{\AgdaPostulate{]ᵐˡ}}\<%
\end{code}

\end{minipage}\hfill
\begin{minipage}[c]{0.42\textwidth}\centering
  \ufig[0.7]{nf4}
\end{minipage}

\medskip\noindent
\emph{The table.}  The coset table is the paper's ``push a gate through
a normal box'' as a structurally recursive function \aid{ract}
(5\,600 lines), with its soundness law \aid{ract-sound}:

%% Generated by make agda from lagda/ex-table.lagda.tex (agda --latex --only-scope-checking);
%% edit that file, not this one.  The hidden block renders nothing.
\begin{code}[hide]%
\>[0]\AgdaKeyword{module}\AgdaSpace{}%
\AgdaModule{vqupit-short.lagda.ex-table}\AgdaSpace{}%
\AgdaKeyword{where}\<%
\\
%
\\[\AgdaEmptyExtraSkip]%
\>[0]\AgdaKeyword{open}\AgdaSpace{}%
\AgdaKeyword{import}\AgdaSpace{}%
\AgdaModule{Data.Product}\AgdaSpace{}%
\AgdaKeyword{using}\AgdaSpace{}%
\AgdaSymbol{(}\AgdaOperator{\AgdaFunction{\AgdaUnderscore{}×\AgdaUnderscore{}}}\AgdaSymbol{;}\AgdaSpace{}%
\AgdaOperator{\AgdaInductiveConstructor{\AgdaUnderscore{},\AgdaUnderscore{}}}\AgdaSymbol{)}\<%
\\
\>[0]\AgdaKeyword{open}\AgdaSpace{}%
\AgdaKeyword{import}\AgdaSpace{}%
\AgdaModule{Word.Base}\AgdaSpace{}%
\AgdaKeyword{using}\AgdaSpace{}%
\AgdaSymbol{(}\AgdaOperator{\AgdaInductiveConstructor{[\AgdaUnderscore{}]ʷ}}\AgdaSymbol{)}\<%
\\
\>[0]\AgdaKeyword{open}\AgdaSpace{}%
\AgdaKeyword{import}\AgdaSpace{}%
\AgdaModule{vqupit-short.lagda.Prelude}\<%
\\
%
\\[\AgdaEmptyExtraSkip]%
\>[0]\AgdaKeyword{private}\AgdaSpace{}%
\AgdaKeyword{variable}\<%
\\
\>[0][@{}l@{\AgdaIndent{0}}]%
\>[2]\AgdaGeneralizable{n}\AgdaSpace{}%
\AgdaSymbol{:}\AgdaSpace{}%
\AgdaDatatype{ℕ}\<%
\\
%
\\[\AgdaEmptyExtraSkip]%
\>[0]\AgdaFunction{Circuit}\AgdaSpace{}%
\AgdaSymbol{:}\AgdaSpace{}%
\AgdaDatatype{ℕ}\AgdaSpace{}%
\AgdaSymbol{→}\AgdaSpace{}%
\AgdaPrimitive{Set}\<%
\\
\>[0]\AgdaFunction{Circuit}\AgdaSpace{}%
\AgdaBound{n}\AgdaSpace{}%
\AgdaSymbol{=}\AgdaSpace{}%
\AgdaDatatype{Word}\AgdaSpace{}%
\AgdaSymbol{(}\AgdaPostulate{Gen}\AgdaSpace{}%
\AgdaBound{n}\AgdaSymbol{)}\<%
\\
%
\\[\AgdaEmptyExtraSkip]%
\>[0]\AgdaKeyword{postulate}\<%
\\
\>[0][@{}l@{\AgdaIndent{0}}]%
\>[2]\AgdaPostulate{ML}\AgdaSpace{}%
\AgdaSymbol{:}\AgdaSpace{}%
\AgdaDatatype{ℕ}\AgdaSpace{}%
\AgdaSymbol{→}\AgdaSpace{}%
\AgdaPrimitive{Set}\<%
\\
%
\>[2]\AgdaOperator{\AgdaPostulate{[\AgdaUnderscore{}]ᶜ}}\AgdaSpace{}%
\AgdaSymbol{:}\AgdaSpace{}%
\AgdaPostulate{ML}\AgdaSpace{}%
\AgdaSymbol{(}\AgdaInductiveConstructor{₁₊}\AgdaSpace{}%
\AgdaGeneralizable{n}\AgdaSymbol{)}\AgdaSpace{}%
\AgdaSymbol{→}\AgdaSpace{}%
\AgdaFunction{Circuit}\AgdaSpace{}%
\AgdaSymbol{(}\AgdaInductiveConstructor{₁₊}\AgdaSpace{}%
\AgdaGeneralizable{n}\AgdaSymbol{)}\<%
\end{code}
\begin{code}%
%
\>[2]\AgdaPostulate{ract}%
\>[14]\AgdaSymbol{:}\AgdaSpace{}%
\AgdaPostulate{ML}\AgdaSpace{}%
\AgdaSymbol{(}\AgdaInductiveConstructor{₁₊}\AgdaSpace{}%
\AgdaGeneralizable{n}\AgdaSymbol{)}\AgdaSpace{}%
\AgdaSymbol{→}\AgdaSpace{}%
\AgdaPostulate{Gen}\AgdaSpace{}%
\AgdaSymbol{(}\AgdaInductiveConstructor{₁₊}\AgdaSpace{}%
\AgdaGeneralizable{n}\AgdaSymbol{)}\AgdaSpace{}%
\AgdaSymbol{→}\AgdaSpace{}%
\AgdaFunction{Circuit}\AgdaSpace{}%
\AgdaGeneralizable{n}\AgdaSpace{}%
\AgdaOperator{\AgdaFunction{×}}\AgdaSpace{}%
\AgdaPostulate{ML}\AgdaSpace{}%
\AgdaSymbol{(}\AgdaInductiveConstructor{₁₊}\AgdaSpace{}%
\AgdaGeneralizable{n}\AgdaSymbol{)}\<%
\\
%
\>[2]\AgdaPostulate{ract-sound}%
\>[14]\AgdaSymbol{:}\AgdaSpace{}%
\AgdaSymbol{∀}\AgdaSpace{}%
\AgdaSymbol{(}\AgdaBound{c}\AgdaSpace{}%
\AgdaSymbol{:}\AgdaSpace{}%
\AgdaPostulate{ML}\AgdaSpace{}%
\AgdaSymbol{(}\AgdaInductiveConstructor{₁₊}\AgdaSpace{}%
\AgdaGeneralizable{n}\AgdaSymbol{))}\AgdaSpace{}%
\AgdaBound{g}\AgdaSpace{}%
\AgdaSymbol{→}\AgdaSpace{}%
\AgdaKeyword{let}\AgdaSpace{}%
\AgdaSymbol{(}\AgdaBound{b'}\AgdaSpace{}%
\AgdaOperator{\AgdaInductiveConstructor{,}}\AgdaSpace{}%
\AgdaBound{c'}\AgdaSymbol{)}\AgdaSpace{}%
\AgdaSymbol{=}\AgdaSpace{}%
\AgdaPostulate{ract}\AgdaSpace{}%
\AgdaBound{c}\AgdaSpace{}%
\AgdaBound{g}\AgdaSpace{}%
\AgdaKeyword{in}\AgdaSpace{}%
\AgdaOperator{\AgdaPostulate{[}}\AgdaSpace{}%
\AgdaBound{c}\AgdaSpace{}%
\AgdaOperator{\AgdaPostulate{]ᶜ}}\AgdaSpace{}%
\AgdaOperator{\AgdaInductiveConstructor{•}}\AgdaSpace{}%
\AgdaOperator{\AgdaInductiveConstructor{[}}\AgdaSpace{}%
\AgdaBound{g}\AgdaSpace{}%
\AgdaOperator{\AgdaInductiveConstructor{]ʷ}}\AgdaSpace{}%
\AgdaOperator{\AgdaPostulate{≈}}\AgdaSpace{}%
\AgdaBound{b'}\AgdaSpace{}%
\AgdaOperator{\AgdaPostulate{↑}}\AgdaSpace{}%
\AgdaOperator{\AgdaInductiveConstructor{•}}\AgdaSpace{}%
\AgdaOperator{\AgdaPostulate{[}}\AgdaSpace{}%
\AgdaBound{c'}\AgdaSpace{}%
\AgdaOperator{\AgdaPostulate{]ᶜ}}\<%
\end{code}


\noindent Each clause of \aid{ract-sound} is derived from the fifteen
simplified symplectic relations, drawn in
Figure~\ref{fig:simplified}; for instance, pushing $S$ through an $E$
box is one $S$-power step, $[\,b\,]^{\mathsf e} \bullet S \;=\;
S^{-b} \bullet S \;\approx\; S^{-(b-1)} \;=\; [\,b-1\,]^{\mathsf
e}$.  A circuit is
rewritten to its normal form by repeating this step, gate by gate;
every rewrite step is thereby traced to a relation.

\emph{Uniqueness.}  The symplectic group acts on
$\aid{Vec}\;\mathrm{Pauli}_1\;n$, a $2n$-dimensional symplectic
$\Zp$-vector space in which $\mathrm{Pauli}_1$ is a two-dimensional
symplectic subspace, represented by a wire in a circuit.  The elements
fixing the bottom wire's subspace pointwise form a copy of
$\mathrm{Sp}(2(n{-}1), \Zp)$, acting above the bottom wire, and the
$\ML$ boxes represent its cosets: every $f \in \Sp$ factors as $f = f'
\circ c_n$ with $c_n \in \ML_n$ and $f'$ above the bottom wire, and
$c_n$ is unique because the bottom-wire output determines the box.
Recursing on $f'$ writes $f$ uniquely as $c_1 \circ c_2 \circ \cdots
\circ c_n$ with $c_k \in \ML_k$, which is Theorem~\ref{thm:unique}.

\medskip\noindent
\emph{Surjectivity.}  Every symplectic transformation is the action of
some circuit, proved constructively by ``clear the bottom wire,
recurse on the tail'' --- the paper's Lemma~3.7, its coset box
computed rather than asserted:

%% Copied from ../vqupit/latex/vqupit/sections/symplectic.tex (symplectic, block 7), the LaTeX
%% that `agda --latex --only-scope-checking' generated for the long version.
\begin{code}%
\>[0]\AgdaFunction{Theorem-LM}\AgdaSpace{}%
\AgdaSymbol{:}\AgdaSpace{}%
\AgdaSymbol{∀}\AgdaSpace{}%
\AgdaSymbol{\{}\AgdaBound{n}\AgdaSymbol{\}}\AgdaSpace{}%
\AgdaSymbol{(}\AgdaBound{p}\AgdaSpace{}%
\AgdaBound{q}\AgdaSpace{}%
\AgdaSymbol{:}\AgdaSpace{}%
\AgdaPostulate{Pauli}\AgdaSpace{}%
\AgdaBound{n}\AgdaSymbol{)}\AgdaSpace{}%
\AgdaSymbol{→}\AgdaSpace{}%
\AgdaPostulate{sform}\AgdaSpace{}%
\AgdaBound{p}\AgdaSpace{}%
\AgdaBound{q}\AgdaSpace{}%
\AgdaOperator{\AgdaDatatype{≡}}\AgdaSpace{}%
\AgdaInductiveConstructor{₁}\AgdaSpace{}%
\AgdaSymbol{→}\<%
\\
\>[0][@{}l@{\AgdaIndent{0}}]%
\>[2]\AgdaFunction{∃}\AgdaSpace{}%
\AgdaSymbol{λ}\AgdaSpace{}%
\AgdaSymbol{(}\AgdaBound{lm}\AgdaSpace{}%
\AgdaSymbol{:}\AgdaSpace{}%
\AgdaFunction{ML}\AgdaSpace{}%
\AgdaBound{n}\AgdaSymbol{)}\AgdaSpace{}%
\AgdaSymbol{→}\AgdaSpace{}%
\AgdaPostulate{act}\AgdaSpace{}%
\AgdaOperator{\AgdaPostulate{[}}\AgdaSpace{}%
\AgdaBound{lm}\AgdaSpace{}%
\AgdaOperator{\AgdaPostulate{]ᵐˡ}}\AgdaSpace{}%
\AgdaBound{p}\AgdaSpace{}%
\AgdaOperator{\AgdaDatatype{≡}}\AgdaSpace{}%
\AgdaPostulate{pZ₀}\AgdaSpace{}%
\AgdaOperator{\AgdaFunction{×}}\AgdaSpace{}%
\AgdaPostulate{act}\AgdaSpace{}%
\AgdaOperator{\AgdaPostulate{[}}\AgdaSpace{}%
\AgdaBound{lm}\AgdaSpace{}%
\AgdaOperator{\AgdaPostulate{]ᵐˡ}}\AgdaSpace{}%
\AgdaBound{q}\AgdaSpace{}%
\AgdaOperator{\AgdaDatatype{≡}}\AgdaSpace{}%
\AgdaPostulate{pX₀}\<%
\end{code}


\medskip\noindent
\emph{Completeness.}  Circuits with the same denotation rewrite
through the table to their normal forms, which denote the same map by
soundness and so are equal by Theorem~\ref{thm:unique}: the circuits
are equivalent.  That argument is the framework's
\aid{by-normalization} (Appendix~\ref{sec:permutations}); with
surjectivity, the simplified relations present the symplectic group:

%% Copied from ../vqupit/latex/vqupit/sections/symplectic.tex (symplectic, the simplified-presentation block), the LaTeX
%% that `agda --latex --only-scope-checking' generated for the long version.
\begin{code}%
\>[0]\AgdaFunction{simplified-presentation}\AgdaSpace{}%
\AgdaSymbol{:}\AgdaSpace{}%
\AgdaSymbol{∀}\AgdaSpace{}%
\AgdaBound{n}\AgdaSpace{}%
\AgdaSymbol{→}\AgdaSpace{}%
\AgdaSymbol{(}\AgdaBound{n}\AgdaSpace{}%
\AgdaOperator{\AgdaPostulate{SRel,\AgdaUnderscore{}===\AgdaUnderscore{}}}\AgdaSymbol{)}\AgdaSpace{}%
\AgdaOperator{\AgdaRecord{IsPresentationOf}}\AgdaSpace{}%
\AgdaSymbol{(}\AgdaPostulate{Sp-group}\AgdaSpace{}%
\AgdaBound{n}\AgdaSymbol{)}\<%
\end{code}


\section{Composing: Paulis, the Semidirect Product, and Scalars}\label{sec:composing}

$\PPauli \cong (\Zp \times \Zp)^n$ is presented over $X$, $Z$ on each
wire by $X^{p} = Z^{p} = \varepsilon$ and $ZX = XZ$; its semantic
group is the phase space that $\Sp$ acts on, and its normal form is
$X^{a}Z^{b}$ per wire.  Gluing this to the symplectic presentation of
\S\ref{sec:symplectic} is the semidirect example of
Appendix~\ref{sec:sn-wreath} with $\Sp$ for $S_n$ and $\PPauli$ for
$(\mathbb{Z}_N)^n$: the presentation has both factors' relations plus
the conjugation relations $h\,n\,h^{-1} = \mathit{conj}\,h\,n$, the
action \aid{conj} spelling out the paper's Lemmas~2.24--2.25.  The
theorem's two conditions --- the action respects the relations of
each factor --- are here proved by semantics: the conjugates of
congruent words have equal denotations by symplectic soundness, hence
are congruent by Pauli \emph{completeness} (\aid{respects-Δ}; the
paper's ``Pauli complete'' rules, Definition~4.7).  Every premise
being a theorem, the semidirect presentation of \aid{Pauli⋊Sp} is
unconditional.

\emph{Plumbing.}  The semidirect presentation has generators $X, Z, H,
S, CZ$ and $3 + 15 + 14$ relations (Pauli, symplectic, conjugation);
the paper's rule sets have $H, S, CZ$ alone.  Three isomorphisms of
presentations close the gap (Figure~\ref{fig:route}), through
\texttt{Simplified-V1} (nineteen rules), \texttt{Paper-V0} (sixteen)
and \texttt{Paper-V1} (Theorem~\ref{thm:main}'s
fifteen): generator reduction ($X$ and $Z$ become words) and relation
reduction (the Pauli and conjugation rules become theorems).
\emph{Semidirect to Simplified-V1} is an isomorphism on \emph{different}
alphabets, $f$ sending $S$ to $R = S \bullet Z^{1/2}$ and $h$ sending
$S$ to $Z^{-1/2} \bullet S$; in particular, $f$ induces the embedding
of the symplectic group into the projective Clifford
group.  \texttt{Simplified-V1}'s nineteen rules are the fifteen simplified symplectic ones
plus $(SH)^3 = \varepsilon$, $H^2SH^2S = SH^2SH^2$ and the two rules
pushing $X$ through $CZ$; \aid{f-well-defined} derives from them the
images under $f$ of the three Pauli and the fourteen conjugation
rules (6\,000 lines).  \emph{Simplified-V1 to Paper-V0} is
the identity on circuits: ten axioms are shared, and V0's \texttt{Lemmas}
(4\,750 lines) derives c10--c15 from its own axioms.
\emph{Paper-V0 to Paper-V1} reaches the paper's own rules.  The two
rule sets differ in how the multiplier is spelled as a word:
\texttt{Paper-V0} spells $M_x = R^{x}HR^{x^{-1}}HR^{x}H$ over $R$;
the paper's~(T1) spells $M_x = \aid{SHS'}\;x^{-1}\;x$ over $S$, with
$\aid{SHS'}\;a\;b = Z^{(b-1)/2}X^{(1-a)/2}S^{b}HS^{a}HS^{b}H$ and
\aid{semi-MR} as $M_gS^{g^2}Z^{(g^2-g)/2} = SM_g$.  \texttt{Paper-V1}
restates the four multiplier rules over \aid{SHS'} and drops one:
$M_1$ is now the word $(SH)^3$, so \aid{M-power} at $k = 0$ already
gives $(SH)^3 \approx \varepsilon$, whereas over $R$ it is $(RH)^3$,
a different word, and the axiom had to stay.  The bridge
$R^{a}HR^{b}HR^{a}H \approx \aid{SHS'}\;a\;b$
(\texttt{Shared.PauliBase}) carries the multiplier rules across,
giving Theorem~\ref{thm:main}.

\emph{Scalars.}  \aid{extension-presentation} extends the projective
presentation --- \texttt{Paper-V0}'s, which has \aid{order-SH} as an
axiom --- by the scalar (Figure~\ref{fig:exact}): add $\omega$ and $\omega^{p} = \varepsilon$,
let $\omega$ commute with every gate, and \emph{correct} each rule by
its scalar.  The correction is $\varepsilon$ except on
\aid{order-SH}, the one rule failing on the nose: $(SH)^3 =
\omega^{-1/8}$, exponent $(p^2{-}1)/8 \equiv -1/8$ in $\Zp$.
The multiplier words now denote $\ket{j} \mapsto \ket{xj}$
only up to a scalar.  The semantic extension is
the Weyl cocycle on the natural semidirect product, $c((P,S),(Q,S'))
= \tfrac12\,\mathrm{sform}(P,\ S \cdot Q)$, with $S \cdot Q$ the
symplectic part acting on the Pauli vector and $\mathrm{sform}$ the
symplectic form of the phase space, the exponent in $PQ =
\omega^{\mathrm{sform}(P,Q)}\,QP$; its cocycle identity is, by
bilinearity, symplectic invariance of \aid{sform}, and the
presented group's cocycle is the pullback of $c$ along
Theorem~\ref{thm:main}'s isomorphism, so the two agree by
construction.  The unitary identification uses Weyl multiplication
\cite[Eq.~(4), Theorem~3]{gross2006hudson} and a \emph{linear} Weil
representation with exact covariance
\cite[Theorem~1.3.2]{gurevich2007geometric}: the operators
$U(e,v,S)=\omega^e W(v)\rho(S)$ obey
\[
 U(e,v,S)U(f,w,T)
 =U\bigl(e+f+\tfrac12\mathrm{sform}(v,Sw),v+Sw,ST\bigr).
\]
This is the Heisenberg--Weil representation
\cite[\S 1.3]{gurevich2010bounds}; conjugation on Weyl operators
determines $S$ and $v$, so it is faithful and its scalars are exactly
$\langle\omega I\rangle$ (at width zero, $e\mapsto\omega^e$).
Adjoining $-I$ gives its direct product with $\mathbb Z_2$ and scalars
$\langle-\omega I\rangle$.  These finite phase choices differ from
the full unitary normalizer, which contains $U(1)$.
Unconditionally:

\begin{code}%
\>[0]\AgdaFunction{presentation-exact}\AgdaSpace{}%
\AgdaSymbol{:}\AgdaSpace{}%
\AgdaSymbol{∀}\AgdaSpace{}%
\AgdaSymbol{(}\AgdaBound{n}\AgdaSpace{}%
\AgdaSymbol{:}\AgdaSpace{}%
\AgdaDatatype{ℕ}\AgdaSymbol{)}\AgdaSpace{}%
\AgdaSymbol{→}\<%
\\
\>[0][@{}l@{\AgdaIndent{0}}]%
\>[2]\AgdaKeyword{let}\AgdaSpace{}%
\AgdaBound{Clifford-group}\AgdaSpace{}%
\AgdaSymbol{=}\AgdaSpace{}%
\AgdaSymbol{λ}\AgdaSpace{}%
\AgdaBound{n}\AgdaSpace{}%
\AgdaSymbol{→}\AgdaSpace{}%
\AgdaPostulate{PE.Presented-group}\AgdaSpace{}%
\AgdaBound{n}\AgdaSpace{}%
\AgdaSymbol{(}\AgdaPostulate{exact-generator-data}\AgdaSpace{}%
\AgdaBound{n}\AgdaSymbol{)}\<%
\\
%
\>[2]\AgdaKeyword{in}\AgdaSpace{}%
\AgdaSymbol{(}\AgdaBound{n}\AgdaSpace{}%
\AgdaOperator{\AgdaPostulate{Exact,\AgdaUnderscore{}===\AgdaUnderscore{}}}\AgdaSymbol{)}\AgdaSpace{}%
\AgdaOperator{\AgdaRecord{IsPresentationOf}}\AgdaSpace{}%
\AgdaSymbol{(}\AgdaBound{Clifford-group}\AgdaSpace{}%
\AgdaBound{n}\AgdaSymbol{)}\<%
\end{code}


\noindent The extension is provably non-split.

\emph{Figure~1 as printed} is over $-\omega$, its multipliers and swap
carrying Legendre symbols and $\lambda_p^2$, so it needs the $-1$ that
the exact group (scalars of odd order) lacks.  Adjoining $-1$ as a
central involution presents the exact group times $\mathbb Z_2$,
\aid{Clifford±}, and Figure~1 is isomorphic to that rule set by
$-\omega \mapsto -1\cdot\omega$, each exact rule's scalar pinned by
reading its lift back: the paper's Theorem~4.10, against our group,
%% Copied from ../vqupit/latex/vqupit/sections/composing.tex (composing, the figure1-presentation block), the LaTeX
%% that `agda --latex --only-scope-checking' generated for the long version.
\begin{code}%
\>[0]\AgdaFunction{figure1-presentation}\AgdaSpace{}%
\AgdaSymbol{:}\AgdaSpace{}%
\AgdaSymbol{∀}\AgdaSpace{}%
\AgdaBound{n}\AgdaSpace{}%
\AgdaSymbol{→}\AgdaSpace{}%
\AgdaSymbol{(}\AgdaBound{n}\AgdaSpace{}%
\AgdaOperator{\AgdaPostulate{F,\AgdaUnderscore{}===\AgdaUnderscore{}}}\AgdaSymbol{)}\AgdaSpace{}%
\AgdaOperator{\AgdaRecord{IsPresentationOf}}\AgdaSpace{}%
\AgdaSymbol{(}\AgdaPostulate{Clifford±}\AgdaSpace{}%
\AgdaBound{n}\AgdaSymbol{)}\<%
\end{code}


\section{Discussion, Related Work and Conclusion}\label{sec:discussion}

Everything above is checked by Agda~2.8 with standard library 2.4 under
\texttt{--safe}, no postulates, no holes.  Not formalised are the \emph{glue facts} --- that the
constructed groups are the groups of unitaries generated by $H, S, CZ$
modulo (or not) phases, the paper's Proposition~2.25 and Theorem~2.26
after Weil; no matrix appears in the development --- the isomorphism
of the two \emph{total} central extensions, and the existence of
primitive roots.  The development is 106\,388 lines in 329 files
(Table~\ref{tab:stats}): framework 15\,000, qupit-specific 60\,000, of
which the symplectic factor is 35\,000 and the composition step 23\,000.  \emph{Lessons:} prove well-definedness by semantics; transport, do not re-prove; make the semantics compute.  A language model helped write the Agda code and drafted this paper from the authors' outline.  

\emph{Related work.}  Presentations of gate groups with hand-checked
normal forms run from Selinger's $n$-qubit Clifford
group~\cite{selinger2015clifford} through~\cite{greylyn2014,bian2021un,li2021on,makary2021realstabilizer,bian2023cliffordt2,bian2023cliffordcs3,bian2023thesis,amy2018cnotdihedral,amy2024toffolihadamard}
to the qutrit~\cite{li2025qutrit} and qupit~\cite{bian2026qudit}
rules; the PROP-based
programme~\cite{clement2023complete,clement2024extensions,clement2024minimal,clement2026realcliffordch,blake2026simpler,blake2026polycontrolled,blake2026phaseaffine},
after Lafont~\cite{lafont2003boolean}, and the complete ZX
calculi~\cite{backens2014zx,jeandel2018zx,booth2022qupitzx,poor2023travaganza}
are pen-and-paper too.  Mechanised completeness is rare ($\Pi$ as a
free symmetric rig groupoid~\cite{choudhury2022symmetries} and its
square-root extension~\cite{carette2024squareroots} are coherence
arguments); verified compilers and
logics~\cite{paykin2017qwire,hietala2021voqc,lehmann2025vyzx,zhou2023coqq,chareton2021qbricks,li2022vqo,li2020proq,peduri2025qbc}
prove soundness of rewrites, never completeness; circuit
optimisers~\cite{li2024quarl,paradis2024synthetiq,kang2023modular} search
over rewrites; mathlib~\cite{mathlib2020,mathlib-doc-presented}, the
AFP~\cite{breitner2010freegroups,kharim2023freegroups} and
MathComp~\cite{gonthier2013oddorder} have no verified coset enumeration
against a given concrete group.

\emph{Next} are a verified matrix semantics closing the chain
to $U_{p^n}(\mathbb{Z}[1/p,\omega])$, the qubit Clifford
group~\cite{selinger2015clifford}, the qutrit rules~\cite{li2025qutrit}
as the case $p = 3$, and submission to the Agda standard library.


%% Acknowledgements are automatically suppressed by the `anonymous' option.
\begin{acks}
We thank Neil J.\ Ross and Yuming Zhao, our coauthors on the rule set
this paper verifies, for discussions throughout.  We have used Claude to
help with coding this project in Agda.  We have used Claude to write this
paper according to our outline; the paper was later checked and revised
by us.
\end{acks}

\bibliographystyle{ACM-Reference-Format}
\bibliography{refs}

\appendix
\section{The Rules, the Normal Boxes, and the Size of the Development}\label{sec:appendix}

\begin{center}
\begin{minipage}[t]{0.42\textwidth}\centering\scriptsize
  \begin{tabular}{@{}r@{\;\,}r@{\,}c@{\,}l@{}}
    order-S & \ufig[0.48]{pv0-order-S-l} & $=$ & \ufig[0.48]{pv0-order-S-r}\\[1pt]
    order-H & \ufig[0.48]{pv0-order-H-l} & $=$ & \ufig[0.48]{pv0-order-H-r}\\[1pt]
    M-power & \ufig[0.48]{pv0-M-power-l} & $=$ & \ufig[0.48]{pv0-M-power-r}\\[1pt]
    semi-MR & \ufig[0.48]{pv0-semi-MR-l} & $=$ & \ufig[0.48]{pv0-semi-MR-r}\\[1pt]
    order-CZ & \ufig[0.48]{pv0-order-CZ-l} & $=$ & \ufig[0.48]{pv0-order-CZ-r}\\[1pt]
    order-Ex & \ufig[0.48]{pv0-order-Ex-l} & $=$ & \ufig[0.48]{pv0-order-Ex-r}\\[1pt]
    order-SH & \ufig[0.48]{pv0-order-SH-l} & $=$ & \ufig[0.48]{pv0-order-SH-r}\\
  \end{tabular}
\end{minipage}\hfill
\begin{minipage}[t]{0.55\textwidth}\centering\scriptsize
  \begin{tabular}{@{}r@{\;\,}l@{}}
    comm-CZ-S↑ & \ufig[0.48]{pv0-comm-CZ-Sup}\\[1pt]
    semi-M↑CZ & \ufig[0.48]{pv0-semi-Mup-CZ}\\[1pt]
    semi-Ex-S↑ & \ufig[0.48]{pv0-semi-Ex-Sup}\\[1pt]
    semi-Ex-H↑ & \ufig[0.48]{pv0-semi-Ex-Hup}\\[1pt]
    comm-HHSHHS & \ufig[0.4]{pv0-comm-HHSHHS-l} $=$ \ufig[0.4]{pv0-comm-HHSHHS-r}\\
  \end{tabular}
\end{minipage}\\[3pt]
\scriptsize
\begin{tabular}{@{}r@{\;\,}l@{\qquad}r@{\;\,}l@{}}
  blake-c12 & \ufig[0.48]{pv0-blake-c12} & yang-baxter & \ufig[0.48]{pv0-yang-baxter}\\[2pt]
  cz-slide & \ufig[0.48]{pv0-cz-slide} & semi-CX↑-CZ↓ & \ufig[0.48]{pv0-semi-CXup-CZdn}\\
\end{tabular}
\captionof{figure}{The sixteen rules of \protect\aid{{\_}CRel,{\_}==={\_}}
(\S\ref{sec:background}) as circuit equations, drawn from the Agda
terms.  Diagrams read in matrix-multiplication order, so a word $w
\bullet v$ is drawn with $v$ on the left.  \protect\aid{blake-c12}
relates two circuits implementing controlled-$Z$;
\protect\aid{yang-baxter} and \protect\aid{cz-slide} are the swap
calculus.  With the scalar restored (\S\ref{sec:composing}) only
\protect\aid{order-SH} changes, to $(SH)^3 = \omega^{-1/8}$.}
\Description{The sixteen circuit equations of the rule set.}
\label{fig:rules}
\end{center}

\begin{center}\footnotesize
  \begin{tabular}{@{}r@{\;\,}r@{\,}c@{\,}l@{\qquad}r@{\;\,}r@{\,}c@{\,}l@{}}
    $A_{a,b}$, $a\neq0$ & \ufig[0.58]{box-A-l} & $=$ & \ufig[0.58]{box-A-r} &
    $D_{a,b}$, $a\neq0$ & \ufig[0.58]{box-D-l} & $=$ & \ufig[0.58]{box-D-r}\\
    $A_{0,b}$, $b\neq0$ & \ufig[0.58]{box-A0-l} & $=$ & \ufig[0.58]{box-A0-r} &
    $D_{0,b}$ & \ufig[0.58]{box-D0-l} & $=$ & \ufig[0.58]{box-D0-r}\\
    $B_{a,b}$, $a\neq0$ & \ufig[0.58]{box-B-l} & $=$ & \ufig[0.58]{box-B-r} &
    $E_{b}$ & \ufig[0.58]{box-E-l} & $=$ & \ufig[0.58]{box-E-r}\\
    $B_{0,b}$ & \ufig[0.58]{box-B0-l} & $=$ & \ufig[0.58]{box-B0-r} & & & & \\
  \end{tabular}
\captionof{figure}{The four kinds of normal box (the qupit paper's Figure~6) and
how each is spelled as a circuit by the section of the normal form
(\S\ref{sec:symplectic}).}
\Description{The seven normal-box shapes and their circuit spellings.}
\label{fig:boxes}
\end{center}

\begin{center}
\resizebox{0.95\textwidth}{!}{%
\begin{tikzpicture}[x=1cm,y=1cm,
  box/.style={draw,rounded corners,align=center,font=\scriptsize,inner sep=2.5pt},
  lab/.style={font=\tiny,align=center,inner sep=1pt},
  >={Stealth[length=4pt]}]
  \node[box] (sp)    at (0,0)     {$\Sp$: full rules (17)\\coset tower, unique NF};
  \node[box] (sim)   at (3.35,0)  {$\Sp$: simplified\\rules (15)};
  \node[box] (sd)    at (6.7,0)   {$\PPauli\rtimes\Sp$: semidirect\\relation on $X,Z,H,S,CZ$};
  \node[box] (v1)    at (10.3,0)  {Simplified-V1\\(19 rules, $H,S,CZ$)};
  \node[box] (v0)    at (13.5,0)  {Paper-V0 (16 rules):\\our Figure~1 mod scalars};
  \node[box] (pauli) at (6.7,-1.5){$\PPauli$: $X^p,\ Z^p,\ XZ{=}ZX$};
  \node[box] (ex)    at (10.3,-1.5){Exact: Paper-V0 $+\ \omega$,\\one correction};
  \node[box] (pv1)   at (13.5,-1.5){Paper-V1 (15 rules):\\Figure~1 mod scalars};
  \draw[->] (sp) -- node[lab,above]{identity iso} (sim);
  \draw[->] (sim) -- node[lab,above]{semidirect thm} (sd);
  \draw[->] (pauli) -- node[lab,right]{semidirect thm} (sd);
  \draw[->] (sd) -- node[lab,above]{generator reduction:\\$X,Z \mapsto$ words} (v1);
  \draw[->] (v1) -- node[lab,above]{identity iso} (v0);
  \draw[->] (v0) -- node[lab,right]{identity iso} (pv1);
  \draw[->] (v0.south west) -- node[lab,below,pos=0.55]{central extension thm} (ex.north east);
\end{tikzpicture}}
\captionof{figure}{The chain of presentations: every arrow is a verified
presentation theorem or isomorphism of presentations; the semantic
group is $\Sp$ on the left, $\PPauli \rtimes \Sp$ from the semidirect
level on, and a central extension of it by $\langle\omega\rangle$ at
``Exact''.}
\Description{Boxes for the presentations, joined by arrows labelled with the theorem that connects them.}
\label{fig:route}
\end{center}

\begin{center}\begin{minipage}{\textwidth}\centering
\captionof{table}{Size of the development (lines of Agda incl.\ comments,
\texttt{wc -l}, September 2026).}
\label{tab:stats}
\small
\begin{tabular}{@{}lrr@{}}
\toprule
Component & Files & Lines \\
\midrule
Framework (\texttt{Word}, \texttt{Presentation}, \texttt{Circuit}, \texttt{Normalization}, \texttt{ForStdlib}) & 54 & 14\,848 \\
Symplectic group (normalisation 8\,352; box relations 18\,161; soundness etc.\ 8\,886) & 118 & 35\,399 \\
Projective Pauli group & 3 & 1\,220 \\
Projective Clifford: semidirect, plumbing, shared calculus & 26 & 17\,996 \\
Clifford with scalars & 11 & 5\,029 \\
\midrule
Other examples in the tree (symmetric, cyclic, wreath, Clifford+$T$, qubit Clifford) & 107 & 27\,552 \\
Not on the chain (\texttt{Simplified-V2}), index, notations & 10 & 4\,344 \\
\midrule
Total & 329 & 106\,388 \\
\bottomrule
\end{tabular}
\end{minipage}
\end{center}

\section{Selected Agda Definitions}\label{app:agda}

The normal-form datatype; the five hypotheses against which the library verifies a coset table
(the general form of the soundness law \aid{ract-sound} of
Appendix~\ref{sec:permutations}), in the generic \texttt{Transfer} module; the
six one-qupit rules of the simplified symplectic rule set
(\S\ref{sec:symplectic}), which with \aid{order-CZ}, \aid{comm-CZ-S↓},
\aid{comm-CZ-S↑} and \aid{selinger-c10}--\aid{c15} make its fifteen; and
the conjugation action of the semidirect presentation
(\S\ref{sec:composing}), the paper's Lemmas~2.22--2.23 as words.  All
three are quoted from the development, as is the doubly inductive
normal-form datatype of \S\ref{sec:symplectic}, shown first.

\begingroup\renewcommand{\AgdaCodeStyle}{\scriptsize}
\noindent\begin{minipage}[t]{0.44\textwidth}
\begin{code}%
\>[0]\AgdaFunction{M}\AgdaSpace{}%
\AgdaFunction{L}\AgdaSpace{}%
\AgdaFunction{ML}\AgdaSpace{}%
\AgdaSymbol{:}\AgdaSpace{}%
\AgdaDatatype{ℕ}\AgdaSpace{}%
\AgdaSymbol{→}\AgdaSpace{}%
\AgdaPrimitive{Set}\<%
\\
\>[0]\AgdaFunction{M}\AgdaSpace{}%
\AgdaNumber{0}\AgdaSpace{}%
\AgdaSymbol{=}\AgdaSpace{}%
\AgdaRecord{⊤}\<%
\\
\>[0]\AgdaFunction{M}\AgdaSpace{}%
\AgdaSymbol{(}\AgdaInductiveConstructor{₁₊}\AgdaSpace{}%
\AgdaBound{n}\AgdaSymbol{)}\AgdaSpace{}%
\AgdaSymbol{=}\AgdaSpace{}%
\AgdaDatatype{Vec}\AgdaSpace{}%
\AgdaPostulate{D}\AgdaSpace{}%
\AgdaBound{n}\AgdaSpace{}%
\AgdaOperator{\AgdaFunction{×}}\AgdaSpace{}%
\AgdaPostulate{E}\<%
\\
\>[0]\AgdaFunction{L}\AgdaSpace{}%
\AgdaNumber{0}\AgdaSpace{}%
\AgdaSymbol{=}\AgdaSpace{}%
\AgdaRecord{⊤}\<%
\\
\>[0]\AgdaFunction{L}\AgdaSpace{}%
\AgdaSymbol{(}\AgdaInductiveConstructor{₁₊}\AgdaSpace{}%
\AgdaBound{n}\AgdaSymbol{)}\AgdaSpace{}%
\AgdaSymbol{=}\AgdaSpace{}%
\AgdaDatatype{Vec}\AgdaSpace{}%
\AgdaPostulate{B}\AgdaSpace{}%
\AgdaBound{n}\AgdaSpace{}%
\AgdaOperator{\AgdaFunction{×}}\AgdaSpace{}%
\AgdaPostulate{A}\<%
\\
\>[0]\AgdaFunction{ML}\AgdaSpace{}%
\AgdaNumber{0}\AgdaSpace{}%
\AgdaSymbol{=}\AgdaSpace{}%
\AgdaRecord{⊤}\<%
\\
\>[0]\AgdaFunction{ML}\AgdaSpace{}%
\AgdaNumber{1}\AgdaSpace{}%
\AgdaSymbol{=}\AgdaSpace{}%
\AgdaFunction{M}\AgdaSpace{}%
\AgdaNumber{1}\AgdaSpace{}%
\AgdaOperator{\AgdaFunction{×}}\AgdaSpace{}%
\AgdaFunction{L}\AgdaSpace{}%
\AgdaNumber{1}\<%
\\
\>[0]\AgdaFunction{ML}\AgdaSpace{}%
\AgdaBound{m}\AgdaSymbol{@(}\AgdaInductiveConstructor{₂₊}\AgdaSpace{}%
\AgdaBound{n}\AgdaSymbol{)}%
\>[350I]\AgdaSymbol{=}\AgdaSpace{}%
\AgdaFunction{M}\AgdaSpace{}%
\AgdaBound{m}\AgdaSpace{}%
\AgdaOperator{\AgdaFunction{×}}\AgdaSpace{}%
\AgdaFunction{L}\AgdaSpace{}%
\AgdaBound{m}\<%
\\
\>[.][@{}l@{}]\<[350I]%
\>[12]\AgdaOperator{\AgdaDatatype{⊎}}\AgdaSpace{}%
\AgdaPostulate{D}\AgdaSpace{}%
\AgdaOperator{\AgdaFunction{×}}\AgdaSpace{}%
\AgdaFunction{ML}\AgdaSpace{}%
\AgdaSymbol{(}\AgdaInductiveConstructor{₁₊}\AgdaSpace{}%
\AgdaBound{n}\AgdaSymbol{)}\<%
\end{code}

\end{minipage}\hfill
\begin{minipage}[t]{0.54\textwidth}
\begin{code}%
\>[0]\AgdaFunction{NF}\AgdaSpace{}%
\AgdaSymbol{:}\AgdaSpace{}%
\AgdaDatatype{ℕ}\AgdaSpace{}%
\AgdaSymbol{→}\AgdaSpace{}%
\AgdaPrimitive{Set}\<%
\\
\>[0]\AgdaFunction{NF}\AgdaSpace{}%
\AgdaNumber{0}\AgdaSpace{}%
\AgdaSymbol{=}\AgdaSpace{}%
\AgdaRecord{⊤}\<%
\\
\>[0]\AgdaFunction{NF}\AgdaSpace{}%
\AgdaSymbol{(}\AgdaInductiveConstructor{₁₊}\AgdaSpace{}%
\AgdaBound{n}\AgdaSymbol{)}\AgdaSpace{}%
\AgdaSymbol{=}\AgdaSpace{}%
\AgdaFunction{NF}\AgdaSpace{}%
\AgdaBound{n}\AgdaSpace{}%
\AgdaOperator{\AgdaFunction{×}}\AgdaSpace{}%
\AgdaFunction{ML}\AgdaSpace{}%
\AgdaSymbol{(}\AgdaInductiveConstructor{₁₊}\AgdaSpace{}%
\AgdaBound{n}\AgdaSymbol{)}\<%
\\
%
\\[\AgdaEmptyExtraSkip]%
\>[0]\AgdaOperator{\AgdaFunction{[\AgdaUnderscore{}]}}\AgdaSpace{}%
\AgdaSymbol{:}\AgdaSpace{}%
\AgdaSymbol{∀}\AgdaSpace{}%
\AgdaSymbol{\{}\AgdaBound{n}\AgdaSymbol{\}}\AgdaSpace{}%
\AgdaSymbol{→}\AgdaSpace{}%
\AgdaFunction{NF}\AgdaSpace{}%
\AgdaBound{n}\AgdaSpace{}%
\AgdaSymbol{→}\AgdaSpace{}%
\AgdaDatatype{Word}\AgdaSpace{}%
\AgdaSymbol{(}\AgdaDatatype{Gen}\AgdaSpace{}%
\AgdaBound{n}\AgdaSymbol{)}\<%
\\
\>[0]\AgdaOperator{\AgdaFunction{[\AgdaUnderscore{}]}}\AgdaSpace{}%
\AgdaSymbol{\{}\AgdaNumber{0}\AgdaSymbol{\}}%
\>[12]\AgdaInductiveConstructor{tt}%
\>[23]\AgdaSymbol{=}\AgdaSpace{}%
\AgdaInductiveConstructor{ε}\<%
\\
\>[0]\AgdaOperator{\AgdaFunction{[\AgdaUnderscore{}]}}\AgdaSpace{}%
\AgdaSymbol{\{}\AgdaInductiveConstructor{₁₊}\AgdaSpace{}%
\AgdaBound{n}\AgdaSymbol{\}}%
\>[12]\AgdaSymbol{(}\AgdaBound{nf}\AgdaSpace{}%
\AgdaOperator{\AgdaInductiveConstructor{,}}\AgdaSpace{}%
\AgdaBound{ml}\AgdaSymbol{)}%
\>[23]\AgdaSymbol{=}\AgdaSpace{}%
\AgdaOperator{\AgdaFunction{[}}\AgdaSpace{}%
\AgdaBound{nf}\AgdaSpace{}%
\AgdaOperator{\AgdaFunction{]}}\AgdaSpace{}%
\AgdaOperator{\AgdaPostulate{↑}}\AgdaSpace{}%
\AgdaOperator{\AgdaInductiveConstructor{•}}\AgdaSpace{}%
\AgdaOperator{\AgdaPostulate{[}}\AgdaSpace{}%
\AgdaBound{ml}\AgdaSpace{}%
\AgdaOperator{\AgdaPostulate{]ᵐˡ}}\<%
\end{code}

\end{minipage}
%% Copied from ../vqupit/latex/vqupit/sections/framework.tex (framework, block 2), the LaTeX
%% that `agda --latex --only-scope-checking' generated for the long version.
\begin{code}%
\>[0][@{}l@{\AgdaIndent{1}}]%
\>[2]\AgdaSymbol{(}\AgdaBound{h=⁻¹f-gen}\AgdaSpace{}%
\AgdaSymbol{:}\AgdaSpace{}%
\AgdaSymbol{∀}\AgdaSpace{}%
\AgdaSymbol{(}\AgdaBound{x}\AgdaSpace{}%
\AgdaSymbol{:}\AgdaSpace{}%
\AgdaPostulate{X}\AgdaSymbol{)}\AgdaSpace{}%
\AgdaSymbol{→}\AgdaSpace{}%
\AgdaSymbol{(}\AgdaOperator{\AgdaInductiveConstructor{[}}\AgdaSpace{}%
\AgdaBound{x}\AgdaSpace{}%
\AgdaOperator{\AgdaInductiveConstructor{]ʷ}}\AgdaSpace{}%
\AgdaOperator{\AgdaInductiveConstructor{,}}\AgdaSpace{}%
\AgdaPostulate{I}\AgdaSymbol{)}\AgdaSpace{}%
\AgdaOperator{\AgdaPostulate{\textasciitilde{}}}\AgdaSpace{}%
\AgdaSymbol{((}\AgdaPostulate{h}\AgdaSpace{}%
\AgdaOperator{\AgdaFunction{ᵗ}}\AgdaSymbol{)}\AgdaSpace{}%
\AgdaPostulate{I}\AgdaSpace{}%
\AgdaSymbol{(}\AgdaPostulate{f}\AgdaSpace{}%
\AgdaBound{x}\AgdaSymbol{)))}%
\>[84]\AgdaComment{--\ 1}\<%
\\
%
\>[2]\AgdaSymbol{(}\AgdaBound{h-wd-ax}\AgdaSpace{}%
\AgdaSymbol{:}\AgdaSpace{}%
\AgdaSymbol{∀}\AgdaSpace{}%
\AgdaSymbol{(}\AgdaBound{c}\AgdaSpace{}%
\AgdaSymbol{:}\AgdaSpace{}%
\AgdaPostulate{C}\AgdaSymbol{)\{}\AgdaBound{u}\AgdaSpace{}%
\AgdaBound{t}\AgdaSpace{}%
\AgdaSymbol{:}\AgdaSpace{}%
\AgdaDatatype{Word}\AgdaSpace{}%
\AgdaPostulate{Y}\AgdaSymbol{\}}\AgdaSpace{}%
\AgdaSymbol{→}\AgdaSpace{}%
\AgdaBound{u}\AgdaSpace{}%
\AgdaOperator{\AgdaPostulate{===₂}}\AgdaSpace{}%
\AgdaBound{t}\AgdaSpace{}%
\AgdaSymbol{→}\AgdaSpace{}%
\AgdaSymbol{((}\AgdaPostulate{h}\AgdaSpace{}%
\AgdaOperator{\AgdaFunction{ᵗ}}\AgdaSymbol{)}\AgdaSpace{}%
\AgdaBound{c}\AgdaSpace{}%
\AgdaBound{u}\AgdaSymbol{)}\AgdaSpace{}%
\AgdaOperator{\AgdaPostulate{\textasciitilde{}}}\AgdaSpace{}%
\AgdaSymbol{((}\AgdaPostulate{h}\AgdaSpace{}%
\AgdaOperator{\AgdaFunction{ᵗ}}\AgdaSymbol{)}\AgdaSpace{}%
\AgdaBound{c}\AgdaSpace{}%
\AgdaBound{t}\AgdaSymbol{))}%
\>[84]\AgdaComment{--\ 2}\<%
\\
%
\>[2]\AgdaSymbol{(}\AgdaBound{f-wd-ax}\AgdaSpace{}%
\AgdaSymbol{:}\AgdaSpace{}%
\AgdaSymbol{∀}\AgdaSpace{}%
\AgdaSymbol{\{}\AgdaBound{w}\AgdaSpace{}%
\AgdaBound{v}\AgdaSymbol{\}}\AgdaSpace{}%
\AgdaSymbol{→}\AgdaSpace{}%
\AgdaBound{w}\AgdaSpace{}%
\AgdaOperator{\AgdaPostulate{===₁}}\AgdaSpace{}%
\AgdaBound{v}\AgdaSpace{}%
\AgdaSymbol{→}\AgdaSpace{}%
\AgdaSymbol{(}\AgdaPostulate{f}\AgdaSpace{}%
\AgdaOperator{\AgdaFunction{ʷ}}\AgdaSymbol{)}\AgdaSpace{}%
\AgdaBound{w}\AgdaSpace{}%
\AgdaOperator{\AgdaPostulate{≈₂}}\AgdaSpace{}%
\AgdaSymbol{(}\AgdaPostulate{f}\AgdaSpace{}%
\AgdaOperator{\AgdaFunction{ʷ}}\AgdaSymbol{)}\AgdaSpace{}%
\AgdaBound{v}\AgdaSymbol{)}%
\>[84]\AgdaComment{--\ 3}\<%
\\
%
\>[2]\AgdaSymbol{(}\AgdaBound{[I]≈ε}\AgdaSpace{}%
\AgdaSymbol{:}\AgdaSpace{}%
\AgdaOperator{\AgdaPostulate{[}}\AgdaSpace{}%
\AgdaPostulate{I}\AgdaSpace{}%
\AgdaOperator{\AgdaPostulate{]}}\AgdaSpace{}%
\AgdaOperator{\AgdaPostulate{≈₂}}\AgdaSpace{}%
\AgdaInductiveConstructor{ε}\AgdaSymbol{)}%
\>[84]\AgdaComment{--\ 4}\<%
\\
%
\>[2]\AgdaSymbol{(}\AgdaBound{h=ract}\AgdaSpace{}%
\AgdaSymbol{:}\AgdaSpace{}%
\AgdaSymbol{∀}\AgdaSpace{}%
\AgdaBound{c}\AgdaSpace{}%
\AgdaBound{b}\AgdaSpace{}%
\AgdaSymbol{→}\AgdaSpace{}%
\AgdaKeyword{let}\AgdaSpace{}%
\AgdaSymbol{(}\AgdaBound{b'}\AgdaSpace{}%
\AgdaOperator{\AgdaInductiveConstructor{,}}\AgdaSpace{}%
\AgdaBound{c'}\AgdaSymbol{)}\AgdaSpace{}%
\AgdaSymbol{=}\AgdaSpace{}%
\AgdaPostulate{h}\AgdaSpace{}%
\AgdaBound{c}\AgdaSpace{}%
\AgdaBound{b}\AgdaSpace{}%
\AgdaKeyword{in}\AgdaSpace{}%
\AgdaOperator{\AgdaPostulate{[}}\AgdaSpace{}%
\AgdaBound{c}\AgdaSpace{}%
\AgdaOperator{\AgdaPostulate{]}}\AgdaSpace{}%
\AgdaOperator{\AgdaInductiveConstructor{•}}\AgdaSpace{}%
\AgdaOperator{\AgdaInductiveConstructor{[}}\AgdaSpace{}%
\AgdaBound{b}\AgdaSpace{}%
\AgdaOperator{\AgdaInductiveConstructor{]ʷ}}\AgdaSpace{}%
\AgdaOperator{\AgdaPostulate{≈₂}}\AgdaSpace{}%
\AgdaSymbol{(}\AgdaPostulate{f}\AgdaSpace{}%
\AgdaOperator{\AgdaFunction{ʷ}}\AgdaSymbol{)}\AgdaSpace{}%
\AgdaBound{b'}\AgdaSpace{}%
\AgdaOperator{\AgdaInductiveConstructor{•}}\AgdaSpace{}%
\AgdaOperator{\AgdaPostulate{[}}\AgdaSpace{}%
\AgdaBound{c'}\AgdaSpace{}%
\AgdaOperator{\AgdaPostulate{]}}\AgdaSymbol{)}%
\>[84]\AgdaComment{--\ 5}\<%
\end{code}

\noindent\begin{minipage}[t]{0.5\textwidth}
\begin{code}%
\>[0]\AgdaFunction{order-S}\AgdaSpace{}%
\AgdaSymbol{:}%
\>[20]\AgdaSymbol{(}\AgdaInductiveConstructor{₁₊}\AgdaSpace{}%
\AgdaGeneralizable{n}\AgdaSymbol{)}\AgdaSpace{}%
\AgdaOperator{\AgdaPostulate{SRel,}}%
\>[34]\AgdaPostulate{S}\AgdaSpace{}%
\AgdaOperator{\AgdaFunction{\textasciicircum{}}}\AgdaSpace{}%
\AgdaPostulate{p}\AgdaSpace{}%
\AgdaOperator{\AgdaPostulate{===}}\AgdaSpace{}%
\AgdaInductiveConstructor{ε}\<%
\\
\>[0]\AgdaFunction{order-H}\AgdaSpace{}%
\AgdaSymbol{:}%
\>[20]\AgdaSymbol{(}\AgdaInductiveConstructor{₁₊}\AgdaSpace{}%
\AgdaGeneralizable{n}\AgdaSymbol{)}\AgdaSpace{}%
\AgdaOperator{\AgdaPostulate{SRel,}}%
\>[34]\AgdaPostulate{H}\AgdaSpace{}%
\AgdaOperator{\AgdaFunction{\textasciicircum{}}}\AgdaSpace{}%
\AgdaNumber{2}\AgdaSpace{}%
\AgdaOperator{\AgdaPostulate{===}}\AgdaSpace{}%
\AgdaPostulate{M₋₁}\<%
\\
\>[0]\AgdaFunction{M-power}\AgdaSpace{}%
\AgdaSymbol{:}%
\>[14]\AgdaSymbol{∀}\AgdaSpace{}%
\AgdaBound{k}\AgdaSpace{}%
\AgdaSymbol{→}\AgdaSpace{}%
\AgdaSymbol{(}\AgdaInductiveConstructor{₁₊}\AgdaSpace{}%
\AgdaGeneralizable{n}\AgdaSymbol{)}\AgdaSpace{}%
\AgdaOperator{\AgdaPostulate{SRel,}}%
\>[34]\AgdaPostulate{Mg\textasciicircum{}}\AgdaSpace{}%
\AgdaBound{k}\AgdaSpace{}%
\AgdaOperator{\AgdaPostulate{===}}\AgdaSpace{}%
\AgdaFunction{M}\AgdaSpace{}%
\AgdaSymbol{(}\AgdaPostulate{g\textasciicircum{}}\AgdaSpace{}%
\AgdaBound{k}\AgdaSymbol{)}\<%
\\
\>[0]\AgdaFunction{semi-MS}\AgdaSpace{}%
\AgdaSymbol{:}%
\>[20]\AgdaSymbol{(}\AgdaInductiveConstructor{₁₊}\AgdaSpace{}%
\AgdaGeneralizable{n}\AgdaSymbol{)}\AgdaSpace{}%
\AgdaOperator{\AgdaPostulate{SRel,}}%
\>[34]\AgdaPostulate{Mg}\AgdaSpace{}%
\AgdaOperator{\AgdaInductiveConstructor{•}}\AgdaSpace{}%
\AgdaPostulate{S}\AgdaSpace{}%
\AgdaOperator{\AgdaPostulate{===}}\AgdaSpace{}%
\AgdaPostulate{S\textasciicircum{}}\AgdaSpace{}%
\AgdaSymbol{(}\AgdaPostulate{g}\AgdaSpace{}%
\AgdaOperator{\AgdaPostulate{*}}\AgdaSpace{}%
\AgdaPostulate{g}\AgdaSymbol{)}\AgdaSpace{}%
\AgdaOperator{\AgdaInductiveConstructor{•}}\AgdaSpace{}%
\AgdaPostulate{Mg}\<%
\\
\>[0]\AgdaFunction{semi-M↑CZ}\AgdaSpace{}%
\AgdaSymbol{:}%
\>[20]\AgdaSymbol{(}\AgdaInductiveConstructor{₂₊}\AgdaSpace{}%
\AgdaGeneralizable{n}\AgdaSymbol{)}\AgdaSpace{}%
\AgdaOperator{\AgdaPostulate{SRel,}}%
\>[34]\AgdaPostulate{Mg}\AgdaSpace{}%
\AgdaOperator{\AgdaPostulate{↑}}\AgdaSpace{}%
\AgdaOperator{\AgdaInductiveConstructor{•}}\AgdaSpace{}%
\AgdaPostulate{CZ}\AgdaSpace{}%
\AgdaOperator{\AgdaPostulate{===}}\AgdaSpace{}%
\AgdaPostulate{CZ\textasciicircum{}}\AgdaSpace{}%
\AgdaPostulate{g}\AgdaSpace{}%
\AgdaOperator{\AgdaInductiveConstructor{•}}\AgdaSpace{}%
\AgdaPostulate{Mg}\AgdaSpace{}%
\AgdaOperator{\AgdaPostulate{↑}}\<%
\\
\>[0]\AgdaFunction{semi-M↓CZ}\AgdaSpace{}%
\AgdaSymbol{:}%
\>[20]\AgdaSymbol{(}\AgdaInductiveConstructor{₂₊}\AgdaSpace{}%
\AgdaGeneralizable{n}\AgdaSymbol{)}\AgdaSpace{}%
\AgdaOperator{\AgdaPostulate{SRel,}}%
\>[34]\AgdaPostulate{Mg}\AgdaSpace{}%
\AgdaOperator{\AgdaPostulate{↓}}\AgdaSpace{}%
\AgdaOperator{\AgdaInductiveConstructor{•}}\AgdaSpace{}%
\AgdaPostulate{CZ}\AgdaSpace{}%
\AgdaOperator{\AgdaPostulate{===}}\AgdaSpace{}%
\AgdaPostulate{CZ\textasciicircum{}}\AgdaSpace{}%
\AgdaPostulate{g}\AgdaSpace{}%
\AgdaOperator{\AgdaInductiveConstructor{•}}\AgdaSpace{}%
\AgdaPostulate{Mg}\AgdaSpace{}%
\AgdaOperator{\AgdaPostulate{↓}}\<%
\end{code}

\end{minipage}\hfill
\begin{minipage}[t]{0.47\textwidth}
\begin{code}%
\>[0]\AgdaFunction{conj}\AgdaSpace{}%
\AgdaSymbol{:}\AgdaSpace{}%
\AgdaDatatype{Sym.Gen}\AgdaSpace{}%
\AgdaGeneralizable{n}\AgdaSpace{}%
\AgdaSymbol{→}\AgdaSpace{}%
\AgdaDatatype{XZ.Gen}\AgdaSpace{}%
\AgdaGeneralizable{n}\AgdaSpace{}%
\AgdaSymbol{→}\AgdaSpace{}%
\AgdaDatatype{Word}\AgdaSpace{}%
\AgdaSymbol{(}\AgdaDatatype{XZ.Gen}\AgdaSpace{}%
\AgdaGeneralizable{n}\AgdaSymbol{)}\<%
\\
\>[0]\AgdaFunction{conj}\AgdaSpace{}%
\AgdaInductiveConstructor{Sym.H-gen}%
\>[17]\AgdaInductiveConstructor{X-gen}%
\>[33]\AgdaSymbol{=}\AgdaSpace{}%
\AgdaPostulate{Z}\<%
\\
\>[0]\AgdaFunction{conj}\AgdaSpace{}%
\AgdaInductiveConstructor{Sym.H-gen}%
\>[17]\AgdaInductiveConstructor{Z-gen}%
\>[33]\AgdaSymbol{=}\AgdaSpace{}%
\AgdaPostulate{X}\AgdaSpace{}%
\AgdaOperator{\AgdaFunction{\textasciicircum{}}}\AgdaSpace{}%
\AgdaPostulate{p-1}\<%
\\
\>[0]\AgdaFunction{conj}\AgdaSpace{}%
\AgdaInductiveConstructor{Sym.S-gen}%
\>[17]\AgdaInductiveConstructor{X-gen}%
\>[33]\AgdaSymbol{=}\AgdaSpace{}%
\AgdaPostulate{X}\AgdaSpace{}%
\AgdaOperator{\AgdaInductiveConstructor{•}}\AgdaSpace{}%
\AgdaPostulate{Z}\<%
\\
\>[0]\AgdaFunction{conj}\AgdaSpace{}%
\AgdaInductiveConstructor{Sym.S-gen}%
\>[17]\AgdaInductiveConstructor{Z-gen}%
\>[33]\AgdaSymbol{=}\AgdaSpace{}%
\AgdaPostulate{Z}\<%
\\
\>[0]\AgdaFunction{conj}\AgdaSpace{}%
\AgdaInductiveConstructor{Sym.CZ-gen}%
\>[17]\AgdaInductiveConstructor{X-gen}%
\>[33]\AgdaSymbol{=}\AgdaSpace{}%
\AgdaPostulate{X}\AgdaSpace{}%
\AgdaOperator{\AgdaInductiveConstructor{•}}\AgdaSpace{}%
\AgdaPostulate{Z}\AgdaSpace{}%
\AgdaOperator{\AgdaPostulate{↑}}\<%
\\
\>[0]\AgdaFunction{conj}\AgdaSpace{}%
\AgdaInductiveConstructor{Sym.CZ-gen}%
\>[17]\AgdaInductiveConstructor{Z-gen}%
\>[33]\AgdaSymbol{=}\AgdaSpace{}%
\AgdaPostulate{Z}\<%
\\
\>[0]\AgdaFunction{conj}\AgdaSpace{}%
\AgdaInductiveConstructor{Sym.CZ-gen}%
\>[17]\AgdaSymbol{(}\AgdaInductiveConstructor{X-gen}\AgdaSpace{}%
\AgdaOperator{\AgdaInductiveConstructor{↥}}\AgdaSymbol{)}%
\>[33]\AgdaSymbol{=}\AgdaSpace{}%
\AgdaPostulate{X}\AgdaSpace{}%
\AgdaOperator{\AgdaPostulate{↑}}\AgdaSpace{}%
\AgdaOperator{\AgdaInductiveConstructor{•}}\AgdaSpace{}%
\AgdaPostulate{Z}\<%
\\
\>[0]\AgdaFunction{conj}\AgdaSpace{}%
\AgdaInductiveConstructor{Sym.CZ-gen}%
\>[17]\AgdaSymbol{(}\AgdaInductiveConstructor{Z-gen}\AgdaSpace{}%
\AgdaOperator{\AgdaInductiveConstructor{↥}}\AgdaSymbol{)}%
\>[33]\AgdaSymbol{=}\AgdaSpace{}%
\AgdaPostulate{Z}\AgdaSpace{}%
\AgdaOperator{\AgdaPostulate{↑}}\<%
\\
\>[0]\AgdaFunction{conj}\AgdaSpace{}%
\AgdaInductiveConstructor{Sym.H-gen}%
\>[17]\AgdaSymbol{(}\AgdaBound{xz}\AgdaSpace{}%
\AgdaOperator{\AgdaInductiveConstructor{↥}}\AgdaSymbol{)}%
\>[33]\AgdaSymbol{=}\AgdaSpace{}%
\AgdaOperator{\AgdaInductiveConstructor{[}}\AgdaSpace{}%
\AgdaBound{xz}\AgdaSpace{}%
\AgdaOperator{\AgdaInductiveConstructor{]ʷ}}\AgdaSpace{}%
\AgdaOperator{\AgdaPostulate{↑}}\<%
\\
\>[0]\AgdaFunction{conj}\AgdaSpace{}%
\AgdaInductiveConstructor{Sym.S-gen}%
\>[17]\AgdaSymbol{(}\AgdaBound{xz}\AgdaSpace{}%
\AgdaOperator{\AgdaInductiveConstructor{↥}}\AgdaSymbol{)}%
\>[33]\AgdaSymbol{=}\AgdaSpace{}%
\AgdaOperator{\AgdaInductiveConstructor{[}}\AgdaSpace{}%
\AgdaBound{xz}\AgdaSpace{}%
\AgdaOperator{\AgdaInductiveConstructor{]ʷ}}\AgdaSpace{}%
\AgdaOperator{\AgdaPostulate{↑}}\<%
\\
\>[0]\AgdaFunction{conj}\AgdaSpace{}%
\AgdaInductiveConstructor{Sym.CZ-gen}%
\>[17]\AgdaSymbol{(}\AgdaBound{xz}\AgdaSpace{}%
\AgdaOperator{\AgdaInductiveConstructor{↥}}\AgdaSpace{}%
\AgdaOperator{\AgdaInductiveConstructor{↥}}\AgdaSymbol{)}%
\>[33]\AgdaSymbol{=}\AgdaSpace{}%
\AgdaOperator{\AgdaInductiveConstructor{[}}\AgdaSpace{}%
\AgdaBound{xz}\AgdaSpace{}%
\AgdaOperator{\AgdaInductiveConstructor{]ʷ}}\AgdaSpace{}%
\AgdaOperator{\AgdaPostulate{↑}}\AgdaSpace{}%
\AgdaOperator{\AgdaPostulate{↑}}\<%
\\
\>[0]\AgdaFunction{conj}\AgdaSpace{}%
\AgdaSymbol{(}\AgdaBound{sym}\AgdaSpace{}%
\AgdaOperator{\AgdaInductiveConstructor{↥}}\AgdaSymbol{)}%
\>[17]\AgdaInductiveConstructor{X-gen}%
\>[33]\AgdaSymbol{=}\AgdaSpace{}%
\AgdaPostulate{X}\<%
\\
\>[0]\AgdaFunction{conj}\AgdaSpace{}%
\AgdaSymbol{(}\AgdaBound{sym}\AgdaSpace{}%
\AgdaOperator{\AgdaInductiveConstructor{↥}}\AgdaSymbol{)}%
\>[17]\AgdaInductiveConstructor{Z-gen}%
\>[33]\AgdaSymbol{=}\AgdaSpace{}%
\AgdaPostulate{Z}\<%
\\
\>[0]\AgdaFunction{conj}\AgdaSpace{}%
\AgdaSymbol{(}\AgdaBound{sym}\AgdaSpace{}%
\AgdaOperator{\AgdaInductiveConstructor{↥}}\AgdaSymbol{)}%
\>[17]\AgdaSymbol{(}\AgdaBound{xz}\AgdaSpace{}%
\AgdaOperator{\AgdaInductiveConstructor{↥}}\AgdaSymbol{)}%
\>[33]\AgdaSymbol{=}\AgdaSpace{}%
\AgdaFunction{conj}\AgdaSpace{}%
\AgdaBound{sym}\AgdaSpace{}%
\AgdaBound{xz}\AgdaSpace{}%
\AgdaOperator{\AgdaPostulate{↑}}\<%
\end{code}

\end{minipage}
\endgroup

%% Appendix C: Section 2 of ../vqupit ("An Introduction to the Framework"),
%% the Agda-generated LaTeX copied verbatim (see the file's header).
%% Copied from ../vqupit/latex/vqupit/sections/permutations.tex (Section 2 of the
%% long version, "An Introduction to the Framework"), the LaTeX that
%% `agda --latex --only-scope-checking' generated from
%% ../vqupit/sections/permutations.lagda.tex; used here as an appendix.
%% Local changes: "This section" -> "This appendix", and the sentence that
%% cited the long version's section on the product constructions.
\begin{code}[hide]%
\>[0]\AgdaKeyword{module}\AgdaSpace{}%
\AgdaModule{vqupit.sections.permutations}\AgdaSpace{}%
\AgdaKeyword{where}\<%
\\
%
\\[\AgdaEmptyExtraSkip]%
\>[0]\AgdaKeyword{open}\AgdaSpace{}%
\AgdaKeyword{import}\AgdaSpace{}%
\AgdaModule{Data.Nat}\AgdaSpace{}%
\AgdaKeyword{using}\AgdaSpace{}%
\AgdaSymbol{(}\AgdaDatatype{ℕ}\AgdaSpace{}%
\AgdaSymbol{;}\AgdaSpace{}%
\AgdaInductiveConstructor{suc}\AgdaSymbol{)}\<%
\\
\>[0]\AgdaKeyword{open}\AgdaSpace{}%
\AgdaKeyword{import}\AgdaSpace{}%
\AgdaModule{Data.Product}\AgdaSpace{}%
\AgdaKeyword{using}\AgdaSpace{}%
\AgdaSymbol{(}\AgdaOperator{\AgdaFunction{\AgdaUnderscore{}×\AgdaUnderscore{}}}\AgdaSpace{}%
\AgdaSymbol{;}\AgdaSpace{}%
\AgdaOperator{\AgdaInductiveConstructor{\AgdaUnderscore{},\AgdaUnderscore{}}}\AgdaSpace{}%
\AgdaSymbol{;}\AgdaSpace{}%
\AgdaField{proj₁}\AgdaSpace{}%
\AgdaSymbol{;}\AgdaSpace{}%
\AgdaField{proj₂}\AgdaSymbol{)}\<%
\\
\>[0]\AgdaKeyword{open}\AgdaSpace{}%
\AgdaKeyword{import}\AgdaSpace{}%
\AgdaModule{Data.Unit}\AgdaSpace{}%
\AgdaKeyword{using}\AgdaSpace{}%
\AgdaSymbol{(}\AgdaRecord{⊤}\AgdaSymbol{)}\<%
\\
\>[0]\AgdaKeyword{open}\AgdaSpace{}%
\AgdaKeyword{import}\AgdaSpace{}%
\AgdaModule{Level}\AgdaSpace{}%
\AgdaKeyword{using}\AgdaSpace{}%
\AgdaSymbol{(}\AgdaFunction{0ℓ}\AgdaSpace{}%
\AgdaSymbol{;}\AgdaSpace{}%
\AgdaPostulate{Level}\AgdaSpace{}%
\AgdaSymbol{;}\AgdaSpace{}%
\AgdaOperator{\AgdaPrimitive{\AgdaUnderscore{}⊔\AgdaUnderscore{}}}\AgdaSymbol{)}\<%
\\
\>[0]\AgdaKeyword{open}\AgdaSpace{}%
\AgdaKeyword{import}\AgdaSpace{}%
\AgdaModule{Relation.Binary}\AgdaSpace{}%
\AgdaKeyword{using}\AgdaSpace{}%
\AgdaSymbol{(}\AgdaFunction{Rel}\AgdaSymbol{)}\<%
\\
\>[0]\AgdaKeyword{open}\AgdaSpace{}%
\AgdaKeyword{import}\AgdaSpace{}%
\AgdaModule{Algebra.Bundles}\AgdaSpace{}%
\AgdaKeyword{using}\AgdaSpace{}%
\AgdaSymbol{(}\AgdaRecord{Group}\AgdaSymbol{)}\<%
\\
%
\\[\AgdaEmptyExtraSkip]%
\>[0]\AgdaKeyword{open}\AgdaSpace{}%
\AgdaKeyword{import}\AgdaSpace{}%
\AgdaModule{Notations}\AgdaSpace{}%
\AgdaKeyword{using}\AgdaSpace{}%
\AgdaSymbol{(}\AgdaInductiveConstructor{₁₊}\AgdaSpace{}%
\AgdaSymbol{;}\AgdaSpace{}%
\AgdaInductiveConstructor{₂₊}\AgdaSpace{}%
\AgdaSymbol{;}\AgdaSpace{}%
\AgdaInductiveConstructor{₃₊}\AgdaSymbol{)}\<%
\\
\>[0]\AgdaKeyword{open}\AgdaSpace{}%
\AgdaKeyword{import}\AgdaSpace{}%
\AgdaModule{Word.Base}\AgdaSpace{}%
\AgdaKeyword{using}\AgdaSpace{}%
\AgdaSymbol{(}\AgdaDatatype{Word}\AgdaSpace{}%
\AgdaSymbol{;}\AgdaSpace{}%
\AgdaFunction{WRel}\AgdaSpace{}%
\AgdaSymbol{;}\AgdaSpace{}%
\AgdaOperator{\AgdaInductiveConstructor{[\AgdaUnderscore{}]ʷ}}\AgdaSpace{}%
\AgdaSymbol{;}\AgdaSpace{}%
\AgdaInductiveConstructor{ε}\AgdaSpace{}%
\AgdaSymbol{;}\AgdaSpace{}%
\AgdaOperator{\AgdaInductiveConstructor{\AgdaUnderscore{}•\AgdaUnderscore{}}}\AgdaSpace{}%
\AgdaSymbol{;}\AgdaSpace{}%
\AgdaFunction{wmap}\AgdaSymbol{)}\<%
\\
\>[0]\AgdaKeyword{open}\AgdaSpace{}%
\AgdaKeyword{import}\AgdaSpace{}%
\AgdaModule{Presentation.GroupLike}\AgdaSpace{}%
\AgdaKeyword{using}\AgdaSpace{}%
\AgdaSymbol{(}\AgdaFunction{Grouplike}\AgdaSymbol{)}\<%
\\
\>[0]\AgdaKeyword{open}\AgdaSpace{}%
\AgdaKeyword{import}\AgdaSpace{}%
\AgdaModule{Examples.Groups.Symmetric.Semantics}\AgdaSpace{}%
\AgdaKeyword{using}\AgdaSpace{}%
\AgdaSymbol{(}\AgdaFunction{Permutation′-group}\AgdaSymbol{)}\<%
\\
\>[0]\AgdaKeyword{import}\AgdaSpace{}%
\AgdaModule{Examples.Groups.Symmetric.Presentation}\AgdaSpace{}%
\AgdaSymbol{as}\AgdaSpace{}%
\AgdaModule{SymPres}\<%
\\
%
\\[\AgdaEmptyExtraSkip]%
\>[0]\AgdaKeyword{private}\AgdaSpace{}%
\AgdaKeyword{variable}\<%
\\
\>[0][@{}l@{\AgdaIndent{0}}]%
\>[2]\AgdaGeneralizable{n}\AgdaSpace{}%
\AgdaSymbol{:}\AgdaSpace{}%
\AgdaDatatype{ℕ}\<%
\\
%
\>[2]\AgdaGeneralizable{a}\AgdaSpace{}%
\AgdaGeneralizable{ℓ}\AgdaSpace{}%
\AgdaSymbol{:}\AgdaSpace{}%
\AgdaPostulate{Level}\<%
\\
%
\>[2]\AgdaGeneralizable{X}\AgdaSpace{}%
\AgdaSymbol{:}\AgdaSpace{}%
\AgdaPrimitive{Set}\<%
\\
%
\\[\AgdaEmptyExtraSkip]%
\>[0]\AgdaKeyword{postulate}\<%
\\
\>[0][@{}l@{\AgdaIndent{0}}]%
\>[2]\AgdaPostulate{proof-elided}\AgdaSpace{}%
\AgdaSymbol{:}\AgdaSpace{}%
\AgdaSymbol{∀}\AgdaSpace{}%
\AgdaSymbol{\{}\AgdaBound{ℓ}\AgdaSymbol{\}}\AgdaSpace{}%
\AgdaSymbol{\{}\AgdaBound{A}\AgdaSpace{}%
\AgdaSymbol{:}\AgdaSpace{}%
\AgdaPrimitive{Set}\AgdaSpace{}%
\AgdaBound{ℓ}\AgdaSymbol{\}}\AgdaSpace{}%
\AgdaSymbol{→}\AgdaSpace{}%
\AgdaBound{A}\<%
\\
%
\>[2]\AgdaPostulate{IsGroupIsomorphism}\AgdaSpace{}%
\AgdaSymbol{:}\AgdaSpace{}%
\AgdaSymbol{\{}\AgdaBound{G}\AgdaSpace{}%
\AgdaSymbol{:}\AgdaSpace{}%
\AgdaRecord{Group}\AgdaSpace{}%
\AgdaGeneralizable{a}\AgdaSpace{}%
\AgdaGeneralizable{ℓ}\AgdaSymbol{\}}\AgdaSpace{}%
\AgdaSymbol{→}\AgdaSpace{}%
\AgdaSymbol{(}\AgdaDatatype{Word}\AgdaSpace{}%
\AgdaGeneralizable{X}\AgdaSpace{}%
\AgdaSymbol{→}\AgdaSpace{}%
\AgdaField{Group.Carrier}\AgdaSpace{}%
\AgdaBound{G}\AgdaSymbol{)}\AgdaSpace{}%
\AgdaSymbol{→}\AgdaSpace{}%
\AgdaPrimitive{Set}\AgdaSpace{}%
\AgdaSymbol{(}\AgdaGeneralizable{a}\AgdaSpace{}%
\AgdaOperator{\AgdaPrimitive{⊔}}\AgdaSpace{}%
\AgdaGeneralizable{ℓ}\AgdaSymbol{)}\<%
\end{code}
\section{An Introduction to the Framework}\label{sec:permutations}

The formalisation is built on a library for presentations, normal forms
and completeness of circuit relation
sets~\cite{bian2026framework,qupit-code}.  This appendix introduces it on
its smallest complete case study: circuits built from a single two-wire
gate, the swap $\sigma$, which present the symmetric groups.  The
walkthrough exhibits the pipeline --- gates,
relations, cosets, tower, semantics, uniqueness, completeness,
presentation --- that \S\ref{sec:symplectic} instantiates again for the
symplectic factor, concretely enough that the latter can be read as the
same pipeline with richer gates and a richer action;
\S\ref{sec:composing} glues such presentations together by the
library's product constructions.

\subsection{Gates, generators, and circuits}\label{sec:sn-gates}

A gate set is a family $\aid{Gate} : \mathbb{N} \to \mathsf{Set}$ giving,
for each arity, the gates on that many adjacent wires; for permutations
there is one gate, on two wires.  The generic circuit layer of the
library turns any such family into wire-indexed
\emph{generators}: a gate placed at the bottom of the wires, or a
generator shifted up by one wire.

\begin{code}%
\>[0]\AgdaKeyword{data}\AgdaSpace{}%
\AgdaDatatype{Gate}\AgdaSpace{}%
\AgdaSymbol{:}\AgdaSpace{}%
\AgdaDatatype{ℕ}\AgdaSpace{}%
\AgdaSymbol{→}\AgdaSpace{}%
\AgdaPrimitive{Set}\AgdaSpace{}%
\AgdaKeyword{where}\<%
\\
\>[0][@{}l@{\AgdaIndent{0}}]%
\>[2]\AgdaInductiveConstructor{σ-gate}\AgdaSpace{}%
\AgdaSymbol{:}\AgdaSpace{}%
\AgdaDatatype{Gate}\AgdaSpace{}%
\AgdaNumber{2}\<%
\\
\>[0]\AgdaKeyword{data}\AgdaSpace{}%
\AgdaDatatype{Gen}\AgdaSpace{}%
\AgdaSymbol{:}\AgdaSpace{}%
\AgdaDatatype{ℕ}\AgdaSpace{}%
\AgdaSymbol{→}\AgdaSpace{}%
\AgdaPrimitive{Set}\AgdaSpace{}%
\AgdaKeyword{where}\<%
\\
\>[0][@{}l@{\AgdaIndent{0}}]%
\>[2]\AgdaInductiveConstructor{gate₁}\AgdaSpace{}%
\AgdaSymbol{:}\AgdaSpace{}%
\AgdaDatatype{Gate}\AgdaSpace{}%
\AgdaNumber{1}\AgdaSpace{}%
\AgdaSymbol{→}\AgdaSpace{}%
\AgdaDatatype{Gen}\AgdaSpace{}%
\AgdaSymbol{(}\AgdaInductiveConstructor{₁₊}\AgdaSpace{}%
\AgdaGeneralizable{n}\AgdaSymbol{)}%
\>[34]\AgdaComment{--\ a\ 1-ary\ gate\ on\ the\ bottom\ wire}\<%
\\
%
\>[2]\AgdaInductiveConstructor{gate₂}\AgdaSpace{}%
\AgdaSymbol{:}\AgdaSpace{}%
\AgdaDatatype{Gate}\AgdaSpace{}%
\AgdaNumber{2}\AgdaSpace{}%
\AgdaSymbol{→}\AgdaSpace{}%
\AgdaDatatype{Gen}\AgdaSpace{}%
\AgdaSymbol{(}\AgdaInductiveConstructor{₂₊}\AgdaSpace{}%
\AgdaGeneralizable{n}\AgdaSymbol{)}%
\>[34]\AgdaComment{--\ a\ 2-ary\ gate\ on\ the\ bottom\ two\ wires}\<%
\\
%
\>[2]\AgdaOperator{\AgdaInductiveConstructor{\AgdaUnderscore{}↥}}%
\>[7]\AgdaSymbol{:}\AgdaSpace{}%
\AgdaDatatype{Gen}\AgdaSpace{}%
\AgdaGeneralizable{n}\AgdaSpace{}%
\AgdaSymbol{→}\AgdaSpace{}%
\AgdaDatatype{Gen}\AgdaSpace{}%
\AgdaSymbol{(}\AgdaInductiveConstructor{₁₊}\AgdaSpace{}%
\AgdaGeneralizable{n}\AgdaSymbol{)}%
\>[34]\AgdaComment{--\ the\ same\ generator,\ one\ wire\ higher}\<%
\end{code}
\begin{code}[hide]%
\>[0]\AgdaFunction{Circuit}\AgdaSpace{}%
\AgdaSymbol{:}\AgdaSpace{}%
\AgdaDatatype{ℕ}\AgdaSpace{}%
\AgdaSymbol{→}\AgdaSpace{}%
\AgdaPrimitive{Set}\<%
\\
\>[0]\AgdaFunction{Circuit}\AgdaSpace{}%
\AgdaBound{n}\AgdaSpace{}%
\AgdaSymbol{=}\AgdaSpace{}%
\AgdaDatatype{Word}\AgdaSpace{}%
\AgdaSymbol{(}\AgdaDatatype{Gen}\AgdaSpace{}%
\AgdaBound{n}\AgdaSymbol{)}\<%
\\
%
\\[\AgdaEmptyExtraSkip]%
\>[0]\AgdaFunction{CRel}\AgdaSpace{}%
\AgdaSymbol{:}\AgdaSpace{}%
\AgdaDatatype{ℕ}\AgdaSpace{}%
\AgdaSymbol{→}\AgdaSpace{}%
\AgdaPrimitive{Set₁}\<%
\\
\>[0]\AgdaFunction{CRel}\AgdaSpace{}%
\AgdaBound{n}\AgdaSpace{}%
\AgdaSymbol{=}\AgdaSpace{}%
\AgdaFunction{Rel}\AgdaSpace{}%
\AgdaSymbol{(}\AgdaFunction{Circuit}\AgdaSpace{}%
\AgdaBound{n}\AgdaSymbol{)}\AgdaSpace{}%
\AgdaFunction{0ℓ}\<%
\\
%
\\[\AgdaEmptyExtraSkip]%
\>[0]\AgdaOperator{\AgdaFunction{\AgdaUnderscore{}↑}}\AgdaSpace{}%
\AgdaSymbol{:}\AgdaSpace{}%
\AgdaFunction{Circuit}\AgdaSpace{}%
\AgdaGeneralizable{n}\AgdaSpace{}%
\AgdaSymbol{→}\AgdaSpace{}%
\AgdaFunction{Circuit}\AgdaSpace{}%
\AgdaSymbol{(}\AgdaInductiveConstructor{₁₊}\AgdaSpace{}%
\AgdaGeneralizable{n}\AgdaSymbol{)}\<%
\\
\>[0]\AgdaOperator{\AgdaFunction{\AgdaUnderscore{}↑}}\AgdaSpace{}%
\AgdaSymbol{=}\AgdaSpace{}%
\AgdaFunction{wmap}\AgdaSpace{}%
\AgdaOperator{\AgdaInductiveConstructor{\AgdaUnderscore{}↥}}\<%
\\
%
\\[\AgdaEmptyExtraSkip]%
\>[0]\AgdaKeyword{pattern}\AgdaSpace{}%
\AgdaInductiveConstructor{σ-gen}\AgdaSpace{}%
\AgdaSymbol{=}\AgdaSpace{}%
\AgdaInductiveConstructor{gate₂}\AgdaSpace{}%
\AgdaInductiveConstructor{σ-gate}\<%
\\
%
\\[\AgdaEmptyExtraSkip]%
\>[0]\AgdaFunction{σ}\AgdaSpace{}%
\AgdaSymbol{:}\AgdaSpace{}%
\AgdaFunction{Circuit}\AgdaSpace{}%
\AgdaSymbol{(}\AgdaInductiveConstructor{₂₊}\AgdaSpace{}%
\AgdaGeneralizable{n}\AgdaSymbol{)}\<%
\\
\>[0]\AgdaFunction{σ}\AgdaSpace{}%
\AgdaSymbol{=}\AgdaSpace{}%
\AgdaOperator{\AgdaInductiveConstructor{[}}\AgdaSpace{}%
\AgdaInductiveConstructor{σ-gen}\AgdaSpace{}%
\AgdaOperator{\AgdaInductiveConstructor{]ʷ}}\<%
\\
%
\\[\AgdaEmptyExtraSkip]%
\>[0]\AgdaKeyword{infix}\AgdaSpace{}%
\AgdaNumber{4}\AgdaSpace{}%
\AgdaOperator{\AgdaDatatype{\AgdaUnderscore{}SRel,\AgdaUnderscore{}===\AgdaUnderscore{}}}\AgdaSpace{}%
\AgdaOperator{\AgdaDatatype{\AgdaUnderscore{}VRel,\AgdaUnderscore{}===\AgdaUnderscore{}}}\AgdaSpace{}%
\AgdaOperator{\AgdaPostulate{\AgdaUnderscore{}≈\AgdaUnderscore{}}}\<%
\\
\>[0]\AgdaKeyword{infixr}\AgdaSpace{}%
\AgdaNumber{10}\AgdaSpace{}%
\AgdaOperator{\AgdaInductiveConstructor{σ•\AgdaUnderscore{}}}\<%
\\
%
\\[\AgdaEmptyExtraSkip]%
\>[0]\AgdaKeyword{private}\AgdaSpace{}%
\AgdaKeyword{variable}\<%
\\
\>[0][@{}l@{\AgdaIndent{0}}]%
\>[2]\AgdaGeneralizable{w}\AgdaSpace{}%
\AgdaGeneralizable{v}\AgdaSpace{}%
\AgdaSymbol{:}\AgdaSpace{}%
\AgdaFunction{Circuit}\AgdaSpace{}%
\AgdaGeneralizable{n}\<%
\\
%
\\[\AgdaEmptyExtraSkip]%
\>[0]\AgdaKeyword{postulate}\<%
\\
\>[0][@{}l@{\AgdaIndent{0}}]%
\>[2]\AgdaOperator{\AgdaPostulate{\AgdaUnderscore{}≈\AgdaUnderscore{}}}\AgdaSpace{}%
\AgdaSymbol{:}\AgdaSpace{}%
\AgdaFunction{Circuit}\AgdaSpace{}%
\AgdaGeneralizable{n}\AgdaSpace{}%
\AgdaSymbol{→}\AgdaSpace{}%
\AgdaFunction{Circuit}\AgdaSpace{}%
\AgdaGeneralizable{n}\AgdaSpace{}%
\AgdaSymbol{→}\AgdaSpace{}%
\AgdaPrimitive{Set}\<%
\end{code}

Here $\aid{₁₊}\;n$ and $\aid{₂₊}\;n$ are the library's patterns for
$1 + n$ and $2 + n$, and $n$ is an implicit argument, so
$\aid{gate₂}\;\aid{σ-gate}$ is a generator at any width from two upward.
An $n$-wire \emph{circuit} is a word of generators, $\aid{Circuit}\;n =
\aid{Word}\;(\aid{Gen}\;n)$: a single generator
$[\,g\,]$\aid{ʷ}, the empty circuit \aid{ε} --- $n$ bare wires, the
identity circuit at that width --- or a sequential composition
$c \aidop{•} c'$.  We write $\sigma$ for the one-letter circuit
$[\,\aid{gate₂}\;\aid{σ-gate}\;]$\aid{ʷ} and $c\;{\uparrow}$ for $c$ shifted
up one wire ($\aid{{\_}↑} = \aid{wmap}\;\aid{{\_}↥}$).  Read
diagrammatically, with wires drawn bottom-up and a word drawn left to
right as written, $\sigma$ is a crossing of the bottom two wires, $\sigma\;{\uparrow}$
the same crossing one wire higher, and $\sigma \aidop{•} \sigma\;{\uparrow}
\aidop{•} \sigma$ the three-wire braid on the left of the
\aid{yang-baxter} equation drawn beside its axiom below.

\subsection{Relations: two axioms, structure for free}\label{sec:sn-relations}

The group-specific relations are an inductive family indexed by the wire
count.  Each axiom is stated once, with its gates at the bottom of the
wires and the width left open (the implicit $n$), so that one constructor
covers every width at which the axiom fits; the structural rule
\aid{cong↑} below moves it to any height.  (Agda names may be mixfix:
each underscore in \aid{{\_}SRel,{\_}==={\_}} marks an argument
position, so the family is applied as $n\;\aid{SRel,}\;w \aidrel{===}
v$, read ``$w = v$ is an axiom on $n$ wires''.)

\begin{center}
\begin{minipage}[c]{0.58\linewidth}
\begingroup\setlength{\mathindent}{0pt}
\begin{code}%
\>[0]\AgdaKeyword{data}\AgdaSpace{}%
\AgdaOperator{\AgdaDatatype{\AgdaUnderscore{}SRel,\AgdaUnderscore{}===\AgdaUnderscore{}}}\AgdaSpace{}%
\AgdaSymbol{:}\AgdaSpace{}%
\AgdaSymbol{(}\AgdaBound{n}\AgdaSpace{}%
\AgdaSymbol{:}\AgdaSpace{}%
\AgdaDatatype{ℕ}\AgdaSymbol{)}\AgdaSpace{}%
\AgdaSymbol{→}\AgdaSpace{}%
\AgdaFunction{WRel}\AgdaSpace{}%
\AgdaSymbol{(}\AgdaDatatype{Gen}\AgdaSpace{}%
\AgdaBound{n}\AgdaSymbol{)}\AgdaSpace{}%
\AgdaKeyword{where}\<%
\\
\>[0][@{}l@{\AgdaIndent{0}}]%
\>[2]\AgdaInductiveConstructor{order}%
\>[14]\AgdaSymbol{:}\AgdaSpace{}%
\AgdaSymbol{(}\AgdaInductiveConstructor{₂₊}\AgdaSpace{}%
\AgdaGeneralizable{n}\AgdaSymbol{)}\AgdaSpace{}%
\AgdaOperator{\AgdaDatatype{SRel,}}\AgdaSpace{}%
\AgdaFunction{σ}\AgdaSpace{}%
\AgdaOperator{\AgdaInductiveConstructor{•}}\AgdaSpace{}%
\AgdaFunction{σ}\AgdaSpace{}%
\AgdaOperator{\AgdaDatatype{===}}\AgdaSpace{}%
\AgdaInductiveConstructor{ε}\<%
\\
%
\>[2]\AgdaInductiveConstructor{yang-baxter}\AgdaSpace{}%
\AgdaSymbol{:}\AgdaSpace{}%
\AgdaSymbol{(}\AgdaInductiveConstructor{₃₊}\AgdaSpace{}%
\AgdaGeneralizable{n}\AgdaSymbol{)}\AgdaSpace{}%
\AgdaOperator{\AgdaDatatype{SRel,}}\AgdaSpace{}%
\AgdaFunction{σ}\AgdaSpace{}%
\AgdaOperator{\AgdaInductiveConstructor{•}}\AgdaSpace{}%
\AgdaFunction{σ}\AgdaSpace{}%
\AgdaOperator{\AgdaFunction{↑}}\AgdaSpace{}%
\AgdaOperator{\AgdaInductiveConstructor{•}}\AgdaSpace{}%
\AgdaFunction{σ}\AgdaSpace{}%
\AgdaOperator{\AgdaDatatype{===}}\AgdaSpace{}%
\AgdaFunction{σ}\AgdaSpace{}%
\AgdaOperator{\AgdaFunction{↑}}\AgdaSpace{}%
\AgdaOperator{\AgdaInductiveConstructor{•}}\AgdaSpace{}%
\AgdaFunction{σ}\AgdaSpace{}%
\AgdaOperator{\AgdaInductiveConstructor{•}}\AgdaSpace{}%
\AgdaFunction{σ}\AgdaSpace{}%
\AgdaOperator{\AgdaFunction{↑}}\<%
\end{code}
\endgroup
\end{minipage}\hspace{2.5em}%
\begin{minipage}[c]{0.28\linewidth}\centering
  \ufig[0.35]{sn-order}\\[4pt]
  \ufig[0.35]{sn-yb}
\end{minipage}
\end{center}

\noindent These two axioms are two thirds of the Coxeter presentation of
$S_n$: $\sigma^2 = \varepsilon$ and the braid (Yang--Baxter) relation.
The third Coxeter relation, that swaps on disjoint wires commute, is not
specific to permutations: it is one of the relations every circuit
theory needs, whatever its gates.  By a
\emph{circuit theory} we mean a gate set with gate-specific relations,
presenting for each width $n$ a monoid of $n$-wire circuits under
sequential composition.  Such circuits are the fixed-width slices of a
symmetric monoidal category (a PROP), and the monoidal structure leaves
exactly two traces at fixed width: the shift $c \mapsto c\;{\uparrow}$ is
``tensoring with an identity wire'', so equations must be preserved by
it, and the interchange law says that gates on disjoint wires commute.
(The library also admits $0$-ary gates --- global scalars, available at
every width --- with one more structural rule, that a scalar is
identified with its own shift, \aid{ω↑=ω}; that it then commutes with
every generator is a theorem, once scalars are assumed to commute with
one another, which is vacuous for a single scalar.)  These are the structural rules every
circuit theory shares --- what a
pen-and-paper presentation leaves implicit in ``a congruence compatible
with composition and tensor product'' --- and the circuit layer adds them
generically: \aid{Lift-Relation} extends any
gate-specific family to the full relation set \aid{{\_}VRel,{\_}==={\_}}:

\begin{code}%
\>[0]\AgdaKeyword{data}\AgdaSpace{}%
\AgdaOperator{\AgdaDatatype{\AgdaUnderscore{}VRel,\AgdaUnderscore{}===\AgdaUnderscore{}}}\AgdaSpace{}%
\AgdaSymbol{:}\AgdaSpace{}%
\AgdaSymbol{(}\AgdaBound{n}\AgdaSpace{}%
\AgdaSymbol{:}\AgdaSpace{}%
\AgdaDatatype{ℕ}\AgdaSymbol{)}\AgdaSpace{}%
\AgdaSymbol{→}\AgdaSpace{}%
\AgdaFunction{CRel}\AgdaSpace{}%
\AgdaBound{n}\AgdaSpace{}%
\AgdaKeyword{where}\<%
\\
\>[0][@{}l@{\AgdaIndent{0}}]%
\>[2]\AgdaInductiveConstructor{srel}%
\>[8]\AgdaSymbol{:}\AgdaSpace{}%
\AgdaGeneralizable{n}\AgdaSpace{}%
\AgdaOperator{\AgdaDatatype{SRel,}}\AgdaSpace{}%
\AgdaGeneralizable{w}\AgdaSpace{}%
\AgdaOperator{\AgdaDatatype{===}}\AgdaSpace{}%
\AgdaGeneralizable{v}\AgdaSpace{}%
\AgdaSymbol{→}\AgdaSpace{}%
\AgdaGeneralizable{n}\AgdaSpace{}%
\AgdaOperator{\AgdaDatatype{VRel,}}\AgdaSpace{}%
\AgdaGeneralizable{w}\AgdaSpace{}%
\AgdaOperator{\AgdaDatatype{===}}\AgdaSpace{}%
\AgdaGeneralizable{v}\<%
\\
%
\>[2]\AgdaInductiveConstructor{cong↑}\AgdaSpace{}%
\AgdaSymbol{:}\AgdaSpace{}%
\AgdaGeneralizable{n}\AgdaSpace{}%
\AgdaOperator{\AgdaDatatype{VRel,}}\AgdaSpace{}%
\AgdaGeneralizable{w}\AgdaSpace{}%
\AgdaOperator{\AgdaDatatype{===}}\AgdaSpace{}%
\AgdaGeneralizable{v}\AgdaSpace{}%
\AgdaSymbol{→}\AgdaSpace{}%
\AgdaSymbol{(}\AgdaInductiveConstructor{₁₊}\AgdaSpace{}%
\AgdaGeneralizable{n}\AgdaSymbol{)}\AgdaSpace{}%
\AgdaOperator{\AgdaDatatype{VRel,}}\AgdaSpace{}%
\AgdaGeneralizable{w}\AgdaSpace{}%
\AgdaOperator{\AgdaFunction{↑}}\AgdaSpace{}%
\AgdaOperator{\AgdaDatatype{===}}\AgdaSpace{}%
\AgdaGeneralizable{v}\AgdaSpace{}%
\AgdaOperator{\AgdaFunction{↑}}\<%
\\
%
\>[2]\AgdaInductiveConstructor{comm₁}\AgdaSpace{}%
\AgdaSymbol{:}\AgdaSpace{}%
\AgdaSymbol{(}\AgdaBound{h}\AgdaSpace{}%
\AgdaSymbol{:}\AgdaSpace{}%
\AgdaDatatype{Gate}\AgdaSpace{}%
\AgdaNumber{1}\AgdaSymbol{)}\AgdaSpace{}%
\AgdaSymbol{(}\AgdaBound{g}\AgdaSpace{}%
\AgdaSymbol{:}\AgdaSpace{}%
\AgdaDatatype{Gen}\AgdaSpace{}%
\AgdaGeneralizable{n}\AgdaSymbol{)}\AgdaSpace{}%
\AgdaSymbol{→}\AgdaSpace{}%
\AgdaSymbol{(}\AgdaInductiveConstructor{₁₊}\AgdaSpace{}%
\AgdaGeneralizable{n}\AgdaSymbol{)}\AgdaSpace{}%
\AgdaOperator{\AgdaDatatype{VRel,}}\AgdaSpace{}%
\AgdaOperator{\AgdaInductiveConstructor{[}}\AgdaSpace{}%
\AgdaBound{g}\AgdaSpace{}%
\AgdaOperator{\AgdaInductiveConstructor{↥}}\AgdaSpace{}%
\AgdaOperator{\AgdaInductiveConstructor{]ʷ}}\AgdaSpace{}%
\AgdaOperator{\AgdaInductiveConstructor{•}}\AgdaSpace{}%
\AgdaOperator{\AgdaInductiveConstructor{[}}\AgdaSpace{}%
\AgdaInductiveConstructor{gate₁}\AgdaSpace{}%
\AgdaBound{h}\AgdaSpace{}%
\AgdaOperator{\AgdaInductiveConstructor{]ʷ}}\AgdaSpace{}%
\AgdaOperator{\AgdaDatatype{===}}\AgdaSpace{}%
\AgdaOperator{\AgdaInductiveConstructor{[}}\AgdaSpace{}%
\AgdaInductiveConstructor{gate₁}\AgdaSpace{}%
\AgdaBound{h}\AgdaSpace{}%
\AgdaOperator{\AgdaInductiveConstructor{]ʷ}}\AgdaSpace{}%
\AgdaOperator{\AgdaInductiveConstructor{•}}\AgdaSpace{}%
\AgdaOperator{\AgdaInductiveConstructor{[}}\AgdaSpace{}%
\AgdaBound{g}\AgdaSpace{}%
\AgdaOperator{\AgdaInductiveConstructor{↥}}\AgdaSpace{}%
\AgdaOperator{\AgdaInductiveConstructor{]ʷ}}\<%
\\
%
\>[2]\AgdaInductiveConstructor{comm₂}\AgdaSpace{}%
\AgdaSymbol{:}\AgdaSpace{}%
\AgdaSymbol{(}\AgdaBound{h}\AgdaSpace{}%
\AgdaSymbol{:}\AgdaSpace{}%
\AgdaDatatype{Gate}\AgdaSpace{}%
\AgdaNumber{2}\AgdaSymbol{)}\AgdaSpace{}%
\AgdaSymbol{(}\AgdaBound{g}\AgdaSpace{}%
\AgdaSymbol{:}\AgdaSpace{}%
\AgdaDatatype{Gen}\AgdaSpace{}%
\AgdaGeneralizable{n}\AgdaSymbol{)}\AgdaSpace{}%
\AgdaSymbol{→}\AgdaSpace{}%
\AgdaSymbol{(}\AgdaInductiveConstructor{₂₊}\AgdaSpace{}%
\AgdaGeneralizable{n}\AgdaSymbol{)}\AgdaSpace{}%
\AgdaOperator{\AgdaDatatype{VRel,}}\AgdaSpace{}%
\AgdaOperator{\AgdaInductiveConstructor{[}}\AgdaSpace{}%
\AgdaBound{g}\AgdaSpace{}%
\AgdaOperator{\AgdaInductiveConstructor{↥}}\AgdaSpace{}%
\AgdaOperator{\AgdaInductiveConstructor{↥}}\AgdaSpace{}%
\AgdaOperator{\AgdaInductiveConstructor{]ʷ}}\AgdaSpace{}%
\AgdaOperator{\AgdaInductiveConstructor{•}}\AgdaSpace{}%
\AgdaOperator{\AgdaInductiveConstructor{[}}\AgdaSpace{}%
\AgdaInductiveConstructor{gate₂}\AgdaSpace{}%
\AgdaBound{h}\AgdaSpace{}%
\AgdaOperator{\AgdaInductiveConstructor{]ʷ}}\AgdaSpace{}%
\AgdaOperator{\AgdaDatatype{===}}\AgdaSpace{}%
\AgdaOperator{\AgdaInductiveConstructor{[}}\AgdaSpace{}%
\AgdaInductiveConstructor{gate₂}\AgdaSpace{}%
\AgdaBound{h}\AgdaSpace{}%
\AgdaOperator{\AgdaInductiveConstructor{]ʷ}}\AgdaSpace{}%
\AgdaOperator{\AgdaInductiveConstructor{•}}\AgdaSpace{}%
\AgdaOperator{\AgdaInductiveConstructor{[}}\AgdaSpace{}%
\AgdaBound{g}\AgdaSpace{}%
\AgdaOperator{\AgdaInductiveConstructor{↥}}\AgdaSpace{}%
\AgdaOperator{\AgdaInductiveConstructor{↥}}\AgdaSpace{}%
\AgdaOperator{\AgdaInductiveConstructor{]ʷ}}\<%
\end{code}

From here on,
$\approx$ denotes the monoid congruence generated by
$n\;\aid{VRel,{\_}==={\_}}$: the least relation
containing the axioms and closed under reflexivity, symmetry, transitivity,
congruence of \aidop{•}, associativity, and the unit laws; deciding it is
the word problem.  Because $\sigma$ is its own inverse and shifts of
invertible generators are invertible, the presented monoid is a group; we
witness this by a \aid{Grouplike} proof (every generator has a left
inverse), by induction on generators.

\subsection{Cosets and the staircase normal form}\label{sec:sn-cosets}

The engine behind the normal form is coset enumeration along the chain
$S_1 < S_2 < \dots < S_n < S_{n+1}$, where $S_n$ sits inside $S_{n+1}$ as
the permutations fixing the bottom wire --- syntactically, the circuits
of the form $c\;{\uparrow}$.  Read as permutations, two circuits lie in the
same right coset of $S_n$ exactly when they bring the same wire to the
bottom position; there are $n{+}1$ choices, so $S_n$ has index $n{+}1$,
and a natural representative of each coset is the \emph{descending
staircase} of swaps that carries that wire down.  In the library a coset
is not a set of circuits but a label naming one, an element of the finite type
$\aid{C}\;n$; here the label is the length of the staircase, as a unary
numeral bounded by the width, and the \emph{section} $[\,\_\,]^{\mathsf
c}$ maps it to the staircase circuit.  The representatives on three,
two, and one wires are drawn on the right:

\begin{center}
\begin{minipage}[c]{0.30\linewidth}
\begingroup\setlength{\mathindent}{0pt}
\begin{code}%
\>[0]\AgdaKeyword{data}\AgdaSpace{}%
\AgdaDatatype{C}\AgdaSpace{}%
\AgdaSymbol{:}\AgdaSpace{}%
\AgdaDatatype{ℕ}\AgdaSpace{}%
\AgdaSymbol{→}\AgdaSpace{}%
\AgdaPrimitive{Set}\AgdaSpace{}%
\AgdaKeyword{where}\<%
\\
\>[0][@{}l@{\AgdaIndent{0}}]%
\>[2]\AgdaInductiveConstructor{ε}%
\>[6]\AgdaSymbol{:}\AgdaSpace{}%
\AgdaDatatype{C}\AgdaSpace{}%
\AgdaGeneralizable{n}\<%
\\
%
\>[2]\AgdaOperator{\AgdaInductiveConstructor{σ•\AgdaUnderscore{}}}\AgdaSpace{}%
\AgdaSymbol{:}\AgdaSpace{}%
\AgdaDatatype{C}\AgdaSpace{}%
\AgdaGeneralizable{n}\AgdaSpace{}%
\AgdaSymbol{→}\AgdaSpace{}%
\AgdaDatatype{C}\AgdaSpace{}%
\AgdaSymbol{(}\AgdaInductiveConstructor{₁₊}\AgdaSpace{}%
\AgdaGeneralizable{n}\AgdaSymbol{)}\<%
\\
\>[0]\AgdaOperator{\AgdaFunction{[\AgdaUnderscore{}]ᶜ}}\AgdaSpace{}%
\AgdaSymbol{:}\AgdaSpace{}%
\AgdaDatatype{C}\AgdaSpace{}%
\AgdaGeneralizable{n}\AgdaSpace{}%
\AgdaSymbol{→}\AgdaSpace{}%
\AgdaFunction{Circuit}\AgdaSpace{}%
\AgdaSymbol{(}\AgdaInductiveConstructor{₁₊}\AgdaSpace{}%
\AgdaGeneralizable{n}\AgdaSymbol{)}\<%
\\
\>[0]\AgdaOperator{\AgdaFunction{[}}\AgdaSpace{}%
\AgdaInductiveConstructor{ε}\AgdaSpace{}%
\AgdaOperator{\AgdaFunction{]ᶜ}}%
\>[11]\AgdaSymbol{=}\AgdaSpace{}%
\AgdaInductiveConstructor{ε}\<%
\\
\>[0]\AgdaOperator{\AgdaFunction{[}}\AgdaSpace{}%
\AgdaOperator{\AgdaInductiveConstructor{σ•}}\AgdaSpace{}%
\AgdaBound{c}\AgdaSpace{}%
\AgdaOperator{\AgdaFunction{]ᶜ}}%
\>[11]\AgdaSymbol{=}\AgdaSpace{}%
\AgdaFunction{σ}\AgdaSpace{}%
\AgdaOperator{\AgdaInductiveConstructor{•}}\AgdaSpace{}%
\AgdaSymbol{(}\AgdaOperator{\AgdaFunction{[}}\AgdaSpace{}%
\AgdaBound{c}\AgdaSpace{}%
\AgdaOperator{\AgdaFunction{]ᶜ}}\AgdaSpace{}%
\AgdaOperator{\AgdaFunction{↑}}\AgdaSymbol{)}\<%
\end{code}
\endgroup
\end{minipage}\hspace{2.5em}%
\begin{minipage}[c]{0.42\linewidth}\centering
\begin{tabular}{@{}l@{\;\;}ccc@{}}
  $\aid{C}\;2$: & \ufig[0.35]{c2-0} & \ufig[0.35]{c2-1} & \ufig[0.35]{c2-2}\\
  & $\varepsilon$ & $\sigma{\bullet}\:\varepsilon$ &
    $\sigma{\bullet}\:\sigma{\bullet}\:\varepsilon$\\[4pt]
  $\aid{C}\;1$, $\aid{C}\;0$: & \ufig[0.35]{c1-0} & \ufig[0.35]{c1-1} & \ufig[0.35]{c0-0}\\
  & $\varepsilon$ & $\sigma{\bullet}\:\varepsilon$ & $\varepsilon$
\end{tabular}
\end{minipage}
\end{center}

\noindent So $\aid{C}\;n$ has exactly the $n{+}1$ elements drawn above
for $n = 2, 1, 0$, matching the index of $S_n$ in $S_{n+1}$.  Why this yields
a normal form is the classical argument of computational group
theory~\cite{sims1994computation,holt2005handbook}:
every $f \in S_{n+1}$ lies in one coset, say with representative $c_n$,
so $f = f' \cdot c_n$ with $f' \in S_n$, and recursing on $f'$ writes $f$
uniquely as $c_1 c_2 \cdots c_n$ with $c_k \in \aid{C}\;k$ --- a numeral
in the \emph{factorial number system}, one digit per level, since
$|S_{n+1}| = (n{+}1) \cdot n \cdots 2$.

\emph{Computing} this factorisation is the job of the coset table.
Given a coset, represented by its staircase, and one generator, the
table says where their product lands --- the updated coset --- and what
is left over on the wires above --- the residual circuit --- so that
staircase followed by generator equals residual (shifted up) followed by
updated staircase.
Pushing $\sigma$ into a staircase can lengthen it, cancel with it ($\sigma^2 =
\varepsilon$), or pass it through and shift it up (Yang--Baxter), or
commute past; for a generator that does not touch the bottom wire we can
recurse.  The residual is empty or a single swap here, but in richer
gate sets it can be a longer circuit, which is why the table returns a
circuit on one wire fewer rather than a generator
(Figure~\ref{fig:pushing} draws the four cases):

\begin{center}
\begin{minipage}[c]{0.65\linewidth}
\begingroup\setlength{\mathindent}{0pt}
\begin{code}%
\>[0]\AgdaFunction{ract}\AgdaSpace{}%
\AgdaSymbol{:}\AgdaSpace{}%
\AgdaDatatype{C}\AgdaSpace{}%
\AgdaSymbol{(}\AgdaInductiveConstructor{₁₊}\AgdaSpace{}%
\AgdaGeneralizable{n}\AgdaSymbol{)}\AgdaSpace{}%
\AgdaSymbol{→}\AgdaSpace{}%
\AgdaDatatype{Gen}\AgdaSpace{}%
\AgdaSymbol{(}\AgdaInductiveConstructor{₂₊}\AgdaSpace{}%
\AgdaGeneralizable{n}\AgdaSymbol{)}\AgdaSpace{}%
\AgdaSymbol{→}\AgdaSpace{}%
\AgdaFunction{Circuit}\AgdaSpace{}%
\AgdaSymbol{(}\AgdaInductiveConstructor{₁₊}\AgdaSpace{}%
\AgdaGeneralizable{n}\AgdaSymbol{)}\AgdaSpace{}%
\AgdaOperator{\AgdaFunction{×}}\AgdaSpace{}%
\AgdaDatatype{C}\AgdaSpace{}%
\AgdaSymbol{(}\AgdaInductiveConstructor{₁₊}\AgdaSpace{}%
\AgdaGeneralizable{n}\AgdaSymbol{)}\<%
\\
\>[0]\AgdaFunction{ract}\AgdaSpace{}%
\AgdaInductiveConstructor{ε}\AgdaSpace{}%
\AgdaInductiveConstructor{σ-gen}\AgdaSpace{}%
\AgdaSymbol{=}\AgdaSpace{}%
\AgdaInductiveConstructor{ε}\AgdaSpace{}%
\AgdaOperator{\AgdaInductiveConstructor{,}}\AgdaSpace{}%
\AgdaOperator{\AgdaInductiveConstructor{σ•}}\AgdaSpace{}%
\AgdaInductiveConstructor{ε}%
\>[36]\AgdaComment{--\ grows}\<%
\\
\>[0]\AgdaFunction{ract}\AgdaSpace{}%
\AgdaSymbol{(}\AgdaOperator{\AgdaInductiveConstructor{σ•}}\AgdaSpace{}%
\AgdaInductiveConstructor{ε}\AgdaSymbol{)}\AgdaSpace{}%
\AgdaInductiveConstructor{σ-gen}\AgdaSpace{}%
\AgdaSymbol{=}\AgdaSpace{}%
\AgdaInductiveConstructor{ε}\AgdaSpace{}%
\AgdaOperator{\AgdaInductiveConstructor{,}}\AgdaSpace{}%
\AgdaInductiveConstructor{ε}%
\>[36]\AgdaComment{--\ cancels\ (σ²\ =\ ε)}\<%
\\
\>[0]\AgdaFunction{ract}\AgdaSpace{}%
\AgdaSymbol{(}\AgdaOperator{\AgdaInductiveConstructor{σ•}}\AgdaSpace{}%
\AgdaOperator{\AgdaInductiveConstructor{σ•}}\AgdaSpace{}%
\AgdaBound{c}\AgdaSymbol{)}\AgdaSpace{}%
\AgdaInductiveConstructor{σ-gen}\AgdaSpace{}%
\AgdaSymbol{=}\AgdaSpace{}%
\AgdaFunction{σ}\AgdaSpace{}%
\AgdaOperator{\AgdaInductiveConstructor{,}}\AgdaSpace{}%
\AgdaOperator{\AgdaInductiveConstructor{σ•}}\AgdaSpace{}%
\AgdaOperator{\AgdaInductiveConstructor{σ•}}\AgdaSpace{}%
\AgdaBound{c}%
\>[36]\AgdaComment{--\ passes\ (Yang-Baxter)}\<%
\\
\>[0]\AgdaFunction{ract}\AgdaSpace{}%
\AgdaInductiveConstructor{ε}\AgdaSpace{}%
\AgdaSymbol{(}\AgdaBound{g}\AgdaSpace{}%
\AgdaOperator{\AgdaInductiveConstructor{↥}}\AgdaSymbol{)}\AgdaSpace{}%
\AgdaSymbol{=}\AgdaSpace{}%
\AgdaOperator{\AgdaInductiveConstructor{[}}\AgdaSpace{}%
\AgdaBound{g}\AgdaSpace{}%
\AgdaOperator{\AgdaInductiveConstructor{]ʷ}}\AgdaSpace{}%
\AgdaOperator{\AgdaInductiveConstructor{,}}\AgdaSpace{}%
\AgdaInductiveConstructor{ε}%
\>[36]\AgdaComment{--\ commutes\ past}\<%
\\
\>[0]\AgdaFunction{ract}\AgdaSpace{}%
\AgdaSymbol{(}\AgdaOperator{\AgdaInductiveConstructor{σ•}}\AgdaSpace{}%
\AgdaBound{c}\AgdaSymbol{)}\AgdaSpace{}%
\AgdaSymbol{(}\AgdaBound{g}\AgdaSpace{}%
\AgdaOperator{\AgdaInductiveConstructor{↥}}\AgdaSymbol{)}\AgdaSpace{}%
\AgdaSymbol{=}\AgdaSpace{}%
\AgdaField{proj₁}\AgdaSpace{}%
\AgdaSymbol{(}\AgdaFunction{ract}\AgdaSpace{}%
\AgdaBound{c}\AgdaSpace{}%
\AgdaBound{g}\AgdaSymbol{)}\AgdaSpace{}%
\AgdaOperator{\AgdaFunction{↑}}\AgdaSpace{}%
\AgdaOperator{\AgdaInductiveConstructor{,}}\AgdaSpace{}%
\AgdaOperator{\AgdaInductiveConstructor{σ•}}\AgdaSpace{}%
\AgdaSymbol{(}\AgdaField{proj₂}\AgdaSpace{}%
\AgdaSymbol{(}\AgdaFunction{ract}\AgdaSpace{}%
\AgdaBound{c}\AgdaSpace{}%
\AgdaBound{g}\AgdaSymbol{))}%
\>[64]\AgdaComment{--\ recurses}\<%
\end{code}
\endgroup
\end{minipage}\hspace{0.5em}%
\begin{minipage}[c]{0.30\linewidth}\raggedright
  \ufigT[0.5]{ract-l}\;\scalebox{0.7}{$\approx$}\;\ufigT[0.5]{ract-r}
\end{minipage}
\end{center}

\begin{figure}[h]
\centering
\begin{tabular}{@{}r@{\,}c@{\,}l@{\;\;}l@{\qquad}r@{\,}c@{\,}l@{\;\;}l@{}}
  \ufigT[0.4]{push4-l} & \scalebox{0.7}{$\rightarrow$} & \ufigT[0.4]{push4-r} &
    \small grows &
  \ufigT[0.4]{push3-l} & \scalebox{0.7}{$\rightarrow$} & \ufigT[0.4]{push3-r} &
    \small cancels ($\sigma^2 = \varepsilon$)\\[4pt]
  \ufigT[0.4]{push2-l} & \scalebox{0.7}{$\rightarrow$} & \ufigT[0.4]{push2-r} &
    \small passes (Yang--Baxter) &
  \ufigT[0.4]{push1-l} & \scalebox{0.7}{$\rightarrow$} & \ufigT[0.4]{push1-r} &
    \small commutes past
\end{tabular}
\caption{Pushing a swap through a staircase: the four cases of the coset
table \aid{ract} as circuit rewrites.}
\label{fig:pushing}
\end{figure}

\medskip\noindent
\aid{σ-gen} abbreviates $\aid{gate₂}\;\aid{σ-gate}$.  The last clause is the recursion: a
shifted generator meets a staircase by meeting its tail one wire down.
Each clause is certified by the \emph{soundness law}, pictured on the
right and proved once per table by a short equational derivation from
the axioms:

\begin{code}%
\>[0]\AgdaFunction{ract-sound}\AgdaSpace{}%
\AgdaSymbol{:}\AgdaSpace{}%
\AgdaSymbol{∀}\AgdaSpace{}%
\AgdaBound{c}\AgdaSpace{}%
\AgdaBound{b}\AgdaSpace{}%
\AgdaSymbol{→}\AgdaSpace{}%
\AgdaKeyword{let}\AgdaSpace{}%
\AgdaSymbol{(}\AgdaBound{b'}\AgdaSpace{}%
\AgdaOperator{\AgdaInductiveConstructor{,}}\AgdaSpace{}%
\AgdaBound{c'}\AgdaSymbol{)}\AgdaSpace{}%
\AgdaSymbol{=}\AgdaSpace{}%
\AgdaFunction{ract}\AgdaSpace{}%
\AgdaBound{c}\AgdaSpace{}%
\AgdaBound{b}\AgdaSpace{}%
\AgdaKeyword{in}\AgdaSpace{}%
\AgdaOperator{\AgdaFunction{[}}\AgdaSpace{}%
\AgdaBound{c}\AgdaSpace{}%
\AgdaOperator{\AgdaFunction{]ᶜ}}\AgdaSpace{}%
\AgdaOperator{\AgdaInductiveConstructor{•}}\AgdaSpace{}%
\AgdaOperator{\AgdaInductiveConstructor{[}}\AgdaSpace{}%
\AgdaBound{b}\AgdaSpace{}%
\AgdaOperator{\AgdaInductiveConstructor{]ʷ}}\AgdaSpace{}%
\AgdaOperator{\AgdaPostulate{≈}}\AgdaSpace{}%
\AgdaBound{b'}\AgdaSpace{}%
\AgdaOperator{\AgdaFunction{↑}}\AgdaSpace{}%
\AgdaOperator{\AgdaInductiveConstructor{•}}\AgdaSpace{}%
\AgdaOperator{\AgdaFunction{[}}\AgdaSpace{}%
\AgdaBound{c'}\AgdaSpace{}%
\AgdaOperator{\AgdaFunction{]ᶜ}}\<%
\end{code}
\begin{code}[hide]%
\>[0]\AgdaFunction{ract-sound}\AgdaSpace{}%
\AgdaSymbol{=}\AgdaSpace{}%
\AgdaPostulate{proof-elided}\<%
\end{code}

The stateful word traversal \aid{{\_}ᵗ} extends
\aid{ract} to whole circuits, threading the coset from left to right and
concatenating the residuals; run from the identity coset, it maps an
$(n{+}2)$-wire circuit to an $(n{+}1)$-wire circuit and a coset --- one
level of normalisation.

\emph{Iterating} the level is the coset tower, whose carrier grows by
one digit per level:

\begin{center}
\begin{minipage}[c]{0.56\linewidth}
\begingroup\setlength{\mathindent}{0pt}
\begin{code}%
\>[0]\AgdaFunction{NF}\AgdaSpace{}%
\AgdaSymbol{:}\AgdaSpace{}%
\AgdaDatatype{ℕ}\AgdaSpace{}%
\AgdaSymbol{→}\AgdaSpace{}%
\AgdaPrimitive{Set}\<%
\\
\>[0]\AgdaFunction{NF}\AgdaSpace{}%
\AgdaNumber{0}%
\>[11]\AgdaSymbol{=}\AgdaSpace{}%
\AgdaRecord{⊤}\<%
\\
\>[0]\AgdaFunction{NF}\AgdaSpace{}%
\AgdaSymbol{(}\AgdaInductiveConstructor{suc}\AgdaSpace{}%
\AgdaBound{k}\AgdaSymbol{)}\AgdaSpace{}%
\AgdaSymbol{=}\AgdaSpace{}%
\AgdaFunction{NF}\AgdaSpace{}%
\AgdaBound{k}\AgdaSpace{}%
\AgdaOperator{\AgdaFunction{×}}\AgdaSpace{}%
\AgdaDatatype{C}\AgdaSpace{}%
\AgdaBound{k}%
\>[26]\AgdaComment{--\ ⊤\ ×\ C\ 0\ ×\ C\ 1\ ×\ ⋯\ ×\ C\ k}\<%
\end{code}
\endgroup
\end{minipage}\hspace{2.5em}%
\begin{minipage}[c]{0.16\linewidth}\centering
  \ufig[0.35]{fig3_staircase}
\end{minipage}
\end{center}

\noindent so $\aid{NF}\;n$ has exactly $n!$ elements.  The normal-form
function $\aid{nf} : \aid{Circuit}\;n \to \aid{NF}\;n$ runs the tables;
its \emph{section} $\aid{inv-nf} : \aid{NF}\;n \to \aid{Circuit}\;n$
rebuilds a circuit from its digits by concatenating the staircases,
suitably shifted, as drawn on the right for $(\mathsf{tt}, c_0, c_1,
c_2, c_3) \in \aid{NF}\;4$, the digit $c_0 \in \aid{C}\;0$ being the
empty staircase --- one staircase per level, each shifted up to the
wires of its level: a ``staircase of staircases''.

\subsection{Semantics, uniqueness, and completeness}\label{sec:sn-semantics}

So far everything was syntax.  The development interprets circuits in
the group $\aid{Permutation′-group}\;n$ of permutations of
$\mathsf{Fin}\;n$, with $\sigma$ swapping the bottom two elements.
\emph{Soundness} --- $w \approx v$ implies $\aid{⟦}w\aid{⟧} \doteq
\aid{⟦}v\aid{⟧}$ --- is a small case analysis over the axioms.  The
heart of the matter is the \emph{unique normal form} property, that
normal forms with equal denotations are equal:
$\aid{⟦}\;\aid{inv-nf}\;u\,\aid{⟧} \doteq
\aid{⟦}\;\aid{inv-nf}\;v\,\aid{⟧}$ implies $u \equiv v$.  It is
proved by induction on the tower: the lower digits fix the last point,
so the top digit alone determines where the permutation sends it; equal
denotations therefore force equal top digits, and cancelling them
reduces the claim to the level below.  Soundness and uniqueness are
exactly the premises of the generic lemma \aid{by-normalization}, which
concludes \emph{completeness}: $\aid{⟦}w\aid{⟧} \doteq
\aid{⟦}v\aid{⟧}$ implies $w \approx v$.  Together with surjectivity
(every permutation is a product of adjacent transpositions), these
assemble into the library's \emph{presentation record}: a
\aid{Grouplike} witness, which makes the presented monoid a group, and
a proof that $\aid{⟦\_⟧}$ is a group isomorphism from that group onto
the target (\aid{IsGroupIsomorphism} is the standard library's, between
the group the witness induces and $G$; the record's module plumbing that
says so is elided):

\begin{code}%
\>[0]\AgdaKeyword{record}\AgdaSpace{}%
\AgdaOperator{\AgdaRecord{\AgdaUnderscore{}IsPresentationOf\AgdaUnderscore{}}}\AgdaSpace{}%
\AgdaSymbol{(}\AgdaOperator{\AgdaBound{\AgdaUnderscore{}===\AgdaUnderscore{}}}\AgdaSpace{}%
\AgdaSymbol{:}\AgdaSpace{}%
\AgdaFunction{WRel}\AgdaSpace{}%
\AgdaGeneralizable{X}\AgdaSymbol{)}\AgdaSpace{}%
\AgdaSymbol{(}\AgdaBound{G}\AgdaSpace{}%
\AgdaSymbol{:}\AgdaSpace{}%
\AgdaRecord{Group}\AgdaSpace{}%
\AgdaGeneralizable{a}\AgdaSpace{}%
\AgdaGeneralizable{ℓ}\AgdaSymbol{)}\AgdaSpace{}%
\AgdaSymbol{:}\AgdaSpace{}%
\AgdaPrimitive{Set}\AgdaSpace{}%
\AgdaSymbol{(}\AgdaBound{a}\AgdaSpace{}%
\AgdaOperator{\AgdaPrimitive{⊔}}\AgdaSpace{}%
\AgdaBound{ℓ}\AgdaSymbol{)}\AgdaSpace{}%
\AgdaKeyword{where}\<%
\\
\>[0][@{}l@{\AgdaIndent{0}}]%
\>[2]\AgdaKeyword{field}\<%
\\
\>[2][@{}l@{\AgdaIndent{0}}]%
\>[4]\AgdaField{gl}%
\>[8]\AgdaSymbol{:}\AgdaSpace{}%
\AgdaFunction{Grouplike}\AgdaSpace{}%
\AgdaOperator{\AgdaBound{\AgdaUnderscore{}===\AgdaUnderscore{}}}\<%
\\
%
\>[4]\AgdaOperator{\AgdaField{⟦\AgdaUnderscore{}⟧}}\AgdaSpace{}%
\AgdaSymbol{:}\AgdaSpace{}%
\AgdaDatatype{Word}\AgdaSpace{}%
\AgdaBound{X}\AgdaSpace{}%
\AgdaSymbol{→}\AgdaSpace{}%
\AgdaField{Group.Carrier}\AgdaSpace{}%
\AgdaBound{G}\<%
\\
%
\>[4]\AgdaField{iso}\AgdaSpace{}%
\AgdaSymbol{:}\AgdaSpace{}%
\AgdaPostulate{IsGroupIsomorphism}\AgdaSpace{}%
\AgdaOperator{\AgdaField{⟦\AgdaUnderscore{}⟧}}\<%
\end{code}

\noindent Every presentation theorem in the paper has the shape of this
one:

\begin{code}%
\>[0]\AgdaFunction{symmetric-presentation}\AgdaSpace{}%
\AgdaSymbol{:}\AgdaSpace{}%
\AgdaSymbol{∀}\AgdaSpace{}%
\AgdaBound{n}\AgdaSpace{}%
\AgdaSymbol{→}\AgdaSpace{}%
\AgdaSymbol{(}\AgdaBound{n}\AgdaSpace{}%
\AgdaOperator{\AgdaDatatype{VRel,\AgdaUnderscore{}===\AgdaUnderscore{}}}\AgdaSymbol{)}\AgdaSpace{}%
\AgdaOperator{\AgdaRecord{IsPresentationOf}}\AgdaSpace{}%
\AgdaSymbol{(}\AgdaFunction{Permutation′-group}\AgdaSpace{}%
\AgdaBound{n}\AgdaSymbol{)}\<%
\end{code}
\begin{code}[hide]%
\>[0]\AgdaFunction{symmetric-presentation}\AgdaSpace{}%
\AgdaBound{n}\AgdaSpace{}%
\AgdaSymbol{=}\AgdaSpace{}%
\AgdaFunction{SymPres.presentation}\AgdaSpace{}%
\AgdaSymbol{\{}\AgdaBound{n}\AgdaSymbol{\}}\<%
\end{code}

\noindent read: the monoid of swap circuits modulo \{\aid{order},
\aid{yang-baxter}\} and the structural rules \emph{is} the group of
permutations of $\mathsf{Fin}\;n$.  The symplectic
factor of \S\ref{sec:symplectic} is this pipeline again, with richer
gates and a richer action, and \S\ref{sec:composing} composes such
presentations.


\end{document}
